\documentclass[11pt]{article}
\usepackage[margin=1.1in]{geometry}
\usepackage{amsmath,amssymb,amsthm,mathtools,booktabs,array,tabularx,longtable}
\usepackage[dvipsnames]{xcolor}
\usepackage[colorlinks=true,linkcolor=MidnightBlue,citecolor=MidnightBlue,urlcolor=MidnightBlue]{hyperref}
\usepackage{fancyhdr}
\pagestyle{fancy}\fancyhf{}\renewcommand{\headrulewidth}{0pt}
\fancyfoot[R]{\thepage}
\fancyfoot[L]{\scriptsize\color{gray}Generated by the Transport Lab from approved parts of its store; not a deposited paper. Papers of record: \href{https://zenodo.org/communities/transport-papers}{Zenodo, Transport Papers}}
\fancypagestyle{plain}{\fancyhf{}\fancyfoot[R]{\thepage}\fancyfoot[L]{\scriptsize\color{gray}Generated by the Transport Lab from approved parts of its store; not a deposited paper. Papers of record: \href{https://zenodo.org/communities/transport-papers}{Zenodo, Transport Papers}}\renewcommand{\headrulewidth}{0pt}}
\newtheorem{theorem}{Theorem}[section]
\newtheorem{proposition}[theorem]{Proposition}
\newtheorem{lemma}[theorem]{Lemma}
\newtheorem{corollary}[theorem]{Corollary}
\newtheorem{conjecture}[theorem]{Conjecture}
\newtheorem{refutation}[theorem]{Refutation}
\newtheorem{prediction}[theorem]{Prediction}
\newtheorem{principle}[theorem]{Principle}
\newtheorem{observation}[theorem]{Observation}
\newtheorem{axiom}[theorem]{Axiom}
\theoremstyle{definition}
\newtheorem{definition}[theorem]{Definition}
\newtheorem{fact}[theorem]{Fact}
\newtheorem{assembled}[theorem]{Assembly}
\newtheorem{postulate}[theorem]{Postulate}
\newtheorem{premise}[theorem]{Premise}
\newtheorem{hypothesis}[theorem]{Hypothesis}
\newtheorem{remark}[theorem]{Remark}
\newtheorem{example}[theorem]{Example}
\theoremstyle{plain}
\setlength{\marginparwidth}{0.8in}
\newcommand{\changedmark}[2]{\marginpar{\raggedright\scriptsize\color{#1}#2}}
\providecommand{\Aa}{}\renewcommand{\Aa}{\mathcal{A}}
\providecommand{\Ad}{}\renewcommand{\Ad}{\widetilde{\mathrm{Ad}}}
\providecommand{\Aev}{}\renewcommand{\Aev}{\mathcal{A}^{+}}
\providecommand{\Aod}{}\renewcommand{\Aod}{\mathcal{A}^{-}}
\providecommand{\B}{}\renewcommand{\B}{\mathrm{B}}
\providecommand{\C}{}\renewcommand{\C}{\mathbb{C}}
\providecommand{\CAR}{}\renewcommand{\CAR}{\mathrm{CAR}}
\providecommand{\CS}{}\renewcommand{\CS}{{\mathbb{S}^{+}}}
\providecommand{\CSm}{}\renewcommand{\CSm}{{\mathbb{S}^{-}}}
\providecommand{\Cl}{}\renewcommand{\Cl}{\mathrm{Cl}}
\providecommand{\Con}{}\renewcommand{\Con}{\mathcal{C}}
\providecommand{\Dl}{}\renewcommand{\Dl}{\Delta}
\providecommand{\EN}{}\renewcommand{\EN}{E_{N}}
\providecommand{\End}{}\renewcommand{\End}{\mathrm{End}}
\providecommand{\F}{}\renewcommand{\F}{\mathbb{F}}
\providecommand{\G}{}\renewcommand{\G}{\Gamma}
\providecommand{\HH}{}\renewcommand{\HH}{\mathbb{H}}
\providecommand{\IS}{}\renewcommand{\IS}{I\sigma_3}
\providecommand{\J}{}\renewcommand{\J}{\mathbf{J}}
\providecommand{\LH}{}\renewcommand{\LH}{L_{H}}
\providecommand{\Mat}{}\renewcommand{\Mat}{\mathrm{M}}
\providecommand{\Mo}{}\renewcommand{\Mo}{M^{\circ}}
\providecommand{\Mp}{}\renewcommand{\Mp}{M^{+}}
\providecommand{\Msun}{}\renewcommand{\Msun}{M_{\odot}}
\providecommand{\Mt}{}\renewcommand{\Mt}{\widetilde{M}}
\providecommand{\Ncone}{}\renewcommand{\Ncone}{\mathcal{N}}
\providecommand{\Neff}{}\renewcommand{\Neff}{N_{\mathrm{eff}}}
\providecommand{\PCS}{}\renewcommand{\PCS}{\mathbb{P}{\mathbb{S}^{+}}}
\providecommand{\PT}{}\renewcommand{\PT}{\mathbb{PT}}
\providecommand{\Pin}{}\renewcommand{\Pin}{\mathrm{Pin}}
\providecommand{\Q}{}\renewcommand{\Q}{\mathbb{Q}}
\providecommand{\R}{}\renewcommand{\R}{\mathbb{R}}
\providecommand{\Ric}{}\renewcommand{\Ric}{\mathrm{Ric}}
\providecommand{\SBH}{}\renewcommand{\SBH}{S_{\mathrm{BH}}}
\providecommand{\Sc}{}\renewcommand{\Sc}{\mathbb{S}}
\providecommand{\Sch}{}\renewcommand{\Sch}{\mathrm{P}}
\providecommand{\Sig}{}\renewcommand{\Sig}{\Sigma}
\providecommand{\Som}{}\renewcommand{\Som}{S_\omega}
\providecommand{\Sp}{}\renewcommand{\Sp}{\mathbb{S}}
\providecommand{\Sph}{}\renewcommand{\Sph}{\mathrm{S}}
\providecommand{\Spi}{}\renewcommand{\Spi}{S_\pi}
\providecommand{\Spin}{}\renewcommand{\Spin}{\mathrm{Spin}}
\providecommand{\Sth}{}\renewcommand{\Sth}{S_{\mathrm{th}}}
\providecommand{\T}{}\renewcommand{\T}{\mathbb{T}}
\providecommand{\TH}{}\renewcommand{\TH}{T_{\mathrm{H}}}
\providecommand{\Tr}{}\renewcommand{\Tr}{\operatorname{Tr}}
\providecommand{\UHF}{}\renewcommand{\UHF}{\mathrm{UHF}}
\providecommand{\Z}{}\renewcommand{\Z}{\mathbb{Z}}
\providecommand{\Zc}{}\renewcommand{\Zc}{\mathrm{Z}}
\providecommand{\cmg}{}\renewcommand{\cmg}{\mathrm{cm^{2}\,g^{-1}}}
\providecommand{\dS}{}\renewcommand{\dS}{\mathrm{dS}}
\providecommand{\dsc}{}\renewcommand{\dsc}{\mathrm{dsc}}
\providecommand{\dtPG}{}\renewcommand{\dtPG}{dt}
\providecommand{\dtS}{}\renewcommand{\dtS}{dt_{\mathrm{S}}}
\providecommand{\ellv}{}\renewcommand{\ellv}{\ell}
\providecommand{\er}{}\renewcommand{\er}{e_r}
\providecommand{\et}{}\renewcommand{\et}{e_t}
\providecommand{\ev}{}\renewcommand{\ev}{\mathrm{ev}}
\providecommand{\gl}{}\renewcommand{\gl}{\mathfrak{gl}}
\providecommand{\hatV}{}\renewcommand{\hatV}{\widehat{V}}
\providecommand{\hb}{}\renewcommand{\hb}{\bar{h}}
\providecommand{\hf}{}\renewcommand{\hf}{\hat{f}}
\providecommand{\id}{}\renewcommand{\id}{\mathrm{id}}
\providecommand{\kpc}{}\renewcommand{\kpc}{\mathrm{kpc}}
\providecommand{\lP}{}\renewcommand{\lP}{\ell_{P}}
\providecommand{\mD}{}\renewcommand{\mD}{\mathbf{m}_{D}}
\providecommand{\mb}{}\renewcommand{\mb}{\bar m}
\providecommand{\mnu}{}\renewcommand{\mnu}{\Sigma m_{\nu}}
\providecommand{\nInf}{}\renewcommand{\nInf}{n_\infty}
\providecommand{\nO}{}\renewcommand{\nO}{n_0}
\providecommand{\ninf}{}\renewcommand{\ninf}{n_{\infty}}
\providecommand{\nz}{}\renewcommand{\nz}{n_{0}}
\providecommand{\proj}{}\renewcommand{\proj}{\operatorname{proj}}
\providecommand{\pv}{}\renewcommand{\pv}{\pi}
\providecommand{\rad}{}\renewcommand{\rad}{\mathrm{rad}}
\providecommand{\sgn}{}\renewcommand{\sgn}{\operatorname{sgn}}
\providecommand{\shA}{}\renewcommand{\shA}{\mathcal{A}}
\providecommand{\shC}{}\renewcommand{\shC}{\mathcal{C}}
\providecommand{\shE}{}\renewcommand{\shE}{\mathcal{E}}
\providecommand{\shJ}{}\renewcommand{\shJ}{\mathcal{J}}
\providecommand{\shK}{}\renewcommand{\shK}{\mathcal{K}}
\providecommand{\shL}{}\renewcommand{\shL}{L}
\providecommand{\shW}{}\renewcommand{\shW}{\mathcal{W}}
\providecommand{\sll}{}\renewcommand{\sll}{\mathfrak{sl}}
\providecommand{\so}{}\renewcommand{\so}{\mathfrak{so}}
\providecommand{\su}{}\renewcommand{\su}{\mathfrak{su}}
\providecommand{\supp}{}\renewcommand{\supp}{\operatorname{supp}}
\providecommand{\tr}{}\renewcommand{\tr}{\operatorname{tr}}
\providecommand{\uu}{}\renewcommand{\uu}{\mathfrak{u}}
\providecommand{\br}{}\renewcommand{\br}[1]{\langle #1\rangle}
\providecommand{\bra}{}\renewcommand{\bra}[1]{\langle #1|}
\providecommand{\file}{}\renewcommand{\file}[1]{\texttt{#1}}
\providecommand{\hp}{}\renewcommand{\hp}[2]{( #1 \,|\, #2 )}
\providecommand{\ip}{}\renewcommand{\ip}[2]{\langle #1, #2\rangle}
\providecommand{\ket}{}\renewcommand{\ket}[1]{|#1\rangle}
\providecommand{\rel}{}\renewcommand{\rel}[2]{D\bigl(#1 \,\|\, #2\bigr)}
\providecommand{\rev}{}\renewcommand{\rev}[1]{\widetilde{#1}}
\providecommand{\slashed}{}\renewcommand{\slashed}[1]{{#1}\mkern-10mu/}
\title{\bfseries Einstein's equations in the completion}
\author{Leon Butler\\ \normalsize\href{mailto:lbutler.math@gmail.com}{\texttt{lbutler.math@gmail.com}}\\ \normalsize ORCID: \href{https://orcid.org/0009-0006-4449-5617}{0009-0006-4449-5617}}
\date{Generated by the Transport Lab, 2 November 2026}
\begin{document}
\maketitle
\begin{abstract}
\normalsize
This trail reads Einstein's equations in the completed algebra, where gravity is a pair of gauge fields $\hb$ and $\Omega$ on flat spacetime. Wherever $\hb$ is invertible, the vacuum equations become the constancy of one vector, the scale tractor.

The equations themselves can be traced to horizons: under the local horizon premise and the Clausius relation, Jacobson's argument gives them, with $\Lambda$ an integration constant. The temperature that argument needs follows from a premise with no temperature in it, the null ground state. Its local form is now a theorem for every Hadamard state of the Dirac field.

Two things stay open. The strength of gravity is still an input, matched only by choosing a cut-off. The ground-state clause still needs a time-independence that local conditions do not supply.
\end{abstract}
\newpage
\noindent \textbf{Authorship and disclosure. }The founding idea of the series is the author's: that division by zero in a Clifford algebra should be read not as a prohibition but as a transport into a higher-dimensional algebra, and that the transported division need not be accompanied by an inverse operation. The direction of investigation and the critical review of the results are likewise the author's. The mathematical development and the drafting of the text were carried out in extended collaboration with Claude Fable 5 (Anthropic), an AI system.

\noindent\fbox{\parbox{\dimexpr\linewidth-2\fboxsep-2\fboxrule}{\textbf{This is not a deposited paper.} It was generated by the Transport Lab on 2 November 2026 and assembled from approved parts of its store. It follows results across several papers of the series; each is stated with its proof and credited to the paper that first proved it. The official versions of those papers are on Zenodo, in the Transport Papers community, \url{https://zenodo.org/communities/transport-papers}.}}
\newpage

\section*{Introduction}
The walk has three parts. It first sets up the dictionary with Einstein's theory and follows two routes to his equations, with what each still assumes. It then follows the quantum state on the horizon, which supplies the temperature, and ends by asking where that state could come from.

\section{What we assume}

\begin{postulate}[(C$_{\mathrm{loc}}$): on every local Rindler horizon the entropy of a cut is the integral of a curvature-independent area density (Loc) and, before matter crosses, the modes on the light rays are jointly thermal for the local boost (R$_{\mathrm{loc}}$)]\label{c:synthesis.c-loc-local-horizon-premise}
The series computes a temperature and an entropy for the edges of black holes and of the universe's horizon, using quantum modes it places on those edges. Jacobson showed in 1995 that Einstein's equations of gravity follow if every small piece of horizon around every point, including the horizon seen by an accelerating observer, has a temperature set by the acceleration and an entropy proportional to its area, by the same amount everywhere. This assumption says exactly that, with the quantum modes running continuously along the light rays of each small horizon and being the same field as the matter that falls through it. It has two parts: the entropy of a slice of the horizon is its area times a fixed number that does not depend on how strong the gravitational field is there, and before matter falls through, the modes are thermal at the accelerating observer's temperature. Where there is no gravitational field the second part now follows from a simpler assumption with no temperature in it; the first part is still assumed. With it, the store proves Einstein's equations, with the strength of gravity set by the entropy per area and the cosmological constant left free. The assumption would fail if the entropy per area turned out to depend on the strength of the gravitational field or on matter, or if gravity were found to deviate from Einstein's equations.
\end{postulate}

\begin{postulate}[(FS), the field-space premise of the gravitational action: the varied gravitational connection, a $\Cl(2,4)$ gauge field with nondegenerate soldering, admits a parallel tractor $I$ with $\ip{X}{I} = 1$ (the physical scale), a holonomy reduction to the stabiliser of $I$]\label{c:synthesis.field-space-premise-parallel-frozen-point}
The series borrows a rule, an action, from which the equations of gravity follow, and that rule gives Einstein's equations only if one decides in advance which shapes of the gravitational field are allowed to vary: those that keep one special fixed direction, the frozen point, unchanged everywhere, measured in the actual scale of the spacetime. Without that decision the same rule gives a different theory, conformal gravity. Several attempts to derive the decision from the series' own structures failed; each turned out to be the same assumption in another form. So it is recorded as an assumption of its own. It explains why the cosmological constant is a fixed number, but for the same reason a measurement that dark energy never changes would not count as separate evidence for it. It would be refuted if dark energy were found to change, or if gravity were found to deviate from Einstein's equations where the rule applies.
\end{postulate}

\begin{postulate}[(HD), the Hadamard deposit premise: the descent's deposited state is a Hadamard state on a Cauchy collar of the descent, for the Dirac field and for the scalar fields of the descended phase; the occupation profile over an admissible reference vacuum stays free data]\label{c:synthesis.hadamard-deposit-premise}
Quantum fields in a gravitational field make sense only in well-behaved states, called Hadamard states, which look like empty space at very short distances. The series' account of the universe's earliest moment, the descent, leaves its fields in some state, but nothing in the series yet shows that state is well behaved. This premise assumes it: just after the descent the state is of the well-behaved kind. Since good behaviour, once present, is kept by the equations of motion, every later situation, including black holes formed by collapse, then has well-behaved states, which is what the series' horizon results and the energy of quantum fields need. The premise does not say which well-behaved state it is, how many particles of each energy were left behind, or what pattern of primordial ripples they carry; those remain open inputs.
\end{postulate}

\begin{postulate}[(N$_{+}$): on the continuum of a horizon's null generators the state of the modes is a gauge-invariant quasi-free ground state for the affine null translations along the generators]\label{c:synthesis.n-plus-null-ground-state}
Black holes, the horizon of an expanding universe and the horizon an accelerating observer sees all have a temperature. The series used to assume each temperature with its own premise. This premise replaces them with one assumption that mentions no temperature: the quantum modes running along each horizon's light rays are in their lowest-energy state for sliding along the rays. The store proves that this state, seen from one side of the horizon, is exactly thermal at the expected temperature. The assumption describes a black hole that has always existed in equilibrium; read on the light rays falling into the horizon, its short-distance form also covers a hole formed by collapse, given a simplification of the radiating fields to one direction per light ray. For the horizon an accelerating observer sees, and for the cosmological horizon with fields that keep their form under rescalings of lengths, it amounts to choosing the standard vacuum of the field theory. It still carries an input: that the state is smooth where the two halves of the horizon meet, the condition every derivation of black-hole temperature uses in some form. A black hole that has always existed in equilibrium is an idealisation in a universe with a positive cosmological constant, where no such equilibrium state exists, so the local form of the assumption is the one that matters there; whether the series' states are smooth enough at short distances to give that local form everywhere is still open, though it now comes down to the usual short-distance requirement on states: how such states look at very short distances where gravity is present, assumed before, is now proved.
\end{postulate}

\section{Einstein's equations in the completion}

Each step in this part either translates Einstein's theory into the completed algebra or asks where its equations come from.

We start from the bridge between the two languages.

Gauge Theory Gravity describes gravity by two gauge fields on flat spacetime, $\hb$ and $\Omega$. Wherever $\hb$ is invertible its field equations are equivalent to Einstein's for the derived metric \[ g(a,b) = \hb\^{}{-1}(a)\cdot\hb\^{}{-1}(b) \] and in vacuum to that metric being Ricci-flat.

\begin{fact}[Imported: equivalence with the Einstein equations]\label{c:interior.imported-equivalence-einstein-equations}
Gauge theory gravity rewrites general relativity in the language of geometric algebra. This imported result says that, wherever its basic field can be inverted, its equations say exactly the same as Einstein's equations for the corresponding metric. In empty space this means the metric has vanishing Ricci curvature.
\emph{Source:} A. Lasenby, C. Doran and S. Gull, Gravity, gauge theories and geometric algebra, Phil. Trans. R. Soc. Lond. A 356 (1998), 487--582., \S 4 (\S 4.3, covariant forms of the field equations), App.~C; C. Doran and A. Lasenby, Geometric Algebra for Physicists, Cambridge University Press (2003)., Ch.~13-14.
\end{fact}

To read those equations in the completion we use the tractor bundle: the vectors $\sigma\nInf + \mu + \rho\nO$ of $\Cl(2,4)$, with a natural connection built from the Schouten tensor $\Sch$, a rearranged form of the curvature.

This step records two facts: the connection preserves the square $Q$ of $\Cl(2,4)$, and a Schouten map proportional to the identity forces \[ c = \tfrac{R}{24} \] In particular $\Sch = 0$ exactly when $\Ric = 0$.

\begin{lemma}[Metric parallelism, and the trace of the Schouten condition]\label{c:interior.metric-parallelism-trace-schouten-condition}
This lemma records two elementary facts behind the tractor form of the gravitational field equations. First, the natural derivative on the tractor bundle, which on flat spacetime is a field of rotations of the six-dimensional conformal model, respects that model's own notion of length. Second, if the Schouten tensor, a rearranged form of the gravitational field strength, is a multiple of the metric then the field is Einstein, and the Schouten tensor vanishes exactly when the Ricci curvature does.
\emph{First proved in} T6, Lem 4.2 (\href{https://doi.org/10.5281/zenodo.22721990}{doi:10.5281/zenodo.22721990}).
\end{lemma}
\noindent\emph{Rests on:} \ref{c:interior.imported-tractor-connection}.

With those two facts the vacuum equations can be repackaged.

Carrying the null scale tractor $\nInf$ along with the tractor gauge field gives \[ D\^{}{T}\_{a}\nInf = -\Sch\bigl(\hb\^{}{-1}(a)\bigr) \] So $\nInf$ is parallel exactly when $\Sch = 0$, that is, exactly when the derived metric is Ricci-flat.

\begin{theorem}[Covariant form of the vacuum equations]\label{c:interior.covariant-form-vacuum-equations}
In the interior region, Einstein's equations for empty space can be recast as a single condition on the gravitational fields: a particular marker in the six-dimensional conformal extension of spacetime, attached to the physical scale, must stay constant as it is carried from point to point by a gauge field built from gravity. This result proves that the two formulations are exactly equivalent, and that the marker is a null (lightlike) object. It also notes that merely having some constant marker of this kind is a weaker condition that does not by itself give the field equations.
\emph{First proved in} T6, Thm 4.5 (\href{https://doi.org/10.5281/zenodo.22721990}{doi:10.5281/zenodo.22721990}).
\end{theorem}
\noindent\emph{Rests on:} \ref{c:interior.imported-equivalence-einstein-equations}, \ref{c:interior.imported-tractor-connection}, \ref{c:interior.metric-parallelism-trace-schouten-condition}.

The same rewriting works when a cosmological constant is present.

For an Einstein metric with constant $\Lambda$ the parallel section is \[ I\_g = \nInf + \tfrac{\Lambda}{6}\,\nO \] with squared norm $Q(I_g) = \Lambda/3$. It is null only when $\Lambda = 0$; any nonzero constant moves it off the cone.

\begin{theorem}[$\Lambda \neq 0$: the scale tractor leaves the cone]\label{c:interior.scale-tractor-leaves-cone}
Einstein's equations with a cosmological constant can be written as the statement that one vector of the six-dimensional conformal model, the scale tractor, is unchanged as the gravitational gauge fields carry it from point to point. Its squared length turns out to be the cosmological constant divided by three. With no cosmological constant it has zero length, and any nonzero cosmological constant moves it off the cone of zero-length directions.
\emph{First proved in} T6, Thm 4.7 (\href{https://doi.org/10.5281/zenodo.22721990}{doi:10.5281/zenodo.22721990}).
\end{theorem}
\noindent\emph{Rests on:} \ref{c:dark-energy.imported-frozen-point-geometry-external}.

All of this rewriting would be suspect if it changed the physics, and this step says it does not.

Where the gauge field is invertible and $\Lambda = 0$, the series' description is not an approximation to general relativity but the same theory, its vacuum equations being the parallelism of the scale tractor.

\begin{assembled}[Imported: the exterior is exactly general relativity]\label{c:hawking.imported-exterior-exactly-general-relativity}
Away from the places where the series' new structure matters, its description of gravity is not an approximation to general relativity but the same theory. With no cosmological constant, both the field equations and their geometric repackaging agree exactly with Einstein's theory in that region.
\emph{First proved in} T7, Fact 2.5 (\href{https://doi.org/10.5281/zenodo.22722121}{doi:10.5281/zenodo.22722121}).
\end{assembled}
\noindent\emph{Rests on:} \ref{c:interior.covariant-form-vacuum-equations}, \ref{c:interior.imported-equivalence-einstein-equations}.

So far the equations were taken from Einstein. The series also asks why they hold, and its first answer runs through horizons.

The local horizon premise (C$_{\mathrm{loc}}$) speaks about every local Rindler horizon, the small horizon an accelerating observer sees near any point. Its first clause, (Loc), says the entropy of a cut is $\int \eta\, dA$ with a density $\eta$ that depends neither on the field strength nor on matter, and horizon entropy carries no Euler term. Its second clause, (R$_{\mathrm{loc}}$), says that before matter crosses, the modes on the light rays are thermal for the local boost at rapidity period $2\pi$, which is the temperature $a/2\pi$. It is adopted as a premise, and (Loc) is not derived on either reading of the entropy.

Add the Clausius relation, that heat crossing the patch raises its entropy by heat over temperature, and Jacobson's 1995 argument goes through. The store writes it out with every hypothesis named and obtains \[ $\mathcal{R}$(a) + 8\pi G\_{$\mathrm{cl}$}\,T(a) = (\Lambda + 4\pi G\_{$\mathrm{cl}$}\,\tau)\,a, \qquad G\_{$\mathrm{cl}$} = 1/(4\eta) \] with $\Lambda$ an integration constant. What picks out Einstein's equations over more complicated ones is that $\eta$ ignores the curvature. The route also takes the whole stress as its source, through the null components $T(k,k)$ in every null direction.

\begin{theorem}[Under (C$_{\mathrm{loc}}$) and the Clausius relation, Jacobson's argument gives Einstein's equations with $\Lambda$ an integration constant and $G_{\mathrm{cl}} = 1/(4\eta)$; every further hypothesis named]\label{c:synthesis.einstein-equations-from-c-loc-and-clausius}
In 1995 Ted Jacobson showed that Einstein's equations of gravity follow if every small patch of horizon, as seen by an accelerating observer, has a temperature set by the acceleration and an entropy proportional to its area, and if heat crossing the patch raises its entropy by heat divided by temperature. The series assumes, as a premise, that its own horizon quantum modes give that temperature and that entropy on all such small horizons. This result writes the argument out step by step under that premise and checks each algebraic step exactly. It also lists everything else the argument needs: that the heat is the energy carried by matter and nothing else changes the entropy, that the horizon patch is momentarily neither expanding nor shearing, that matter is described by an ordinary conserved energy tensor, that the gravitational field is smooth enough, and that space-time has four dimensions. Given all that, Einstein's equations follow, with the strength of gravity fixed by the entropy per area and the cosmological constant left free. What singles out Einstein's equations, rather than more complicated versions of them, is that the entropy per area is the same everywhere and does not depend on the strength of the gravitational field there. The derivation is conditional: it shows what the premise would deliver, not that the premise holds.
\end{theorem}
\noindent\emph{Rests on:} \ref{c:synthesis.imported-jacobson-clausius-1995}, \ref{c:synthesis.einstein-condition-is-infinity-vacuum-incidence}, \ref{c:synthesis.imported-raychaudhuri-null-congruence}, \ref{c:synthesis.c-loc-local-horizon-premise}.

The series has a second route to the same equations, through an action.

The borrowed MacDowell-Mansouri functional gives Einstein's equations only if the varied connection admits a parallel tractor $I$ with $\ip{X}{I} = 1$, the frozen point, so that $Q(I) = \Lambda/3$ is constant. Without that restriction the same integrand is the Weyl square, which gives conformal gravity instead. No route derives the restriction from transport, so it is adopted as a premise of its own, the field-space premise (FS), refuted by a confirmed $w \neq -1$.

Kept side by side, the two routes constrain each other. With the action's single coupling $G$ and $\Lambda \neq 0$, the Clausius relation at rapidity period $2\pi c$ holds where matter crosses exactly when $\eta = c/(4G)$. So, given the calibration $\eta = \tfrac14$ at the action's $G$, the one coupling of (Loc) and the period of (R$_{\mathrm{loc}}$) imply each other there, and the Clausius relation becomes a constraint on the thermal data rather than a derivation. In empty regions it says nothing.

\begin{theorem}[If the field equations are those of the field-space action with one normalisation $G$ and $\Lambda \neq 0$, and the Clausius relation holds on local patches at rapidity period $2\pi c$, then at every point where the stress is not a multiple of the metric the Clausius relation is exactly $\eta = c/(4G)$ (Jacobson's identification at period $2\pi c$); so, given the calibration $\eta = \tfrac14$ with the action's $G$ in the same units ($G = 1$), the one-coupling clause of (Loc) and the period of (R$_{\mathrm{loc}}$) imply each other where matter crosses and the stored independence model of each fails there given the other datum, while on matter-free regions the relation is empty]\label{c:synthesis.clausius-on-the-action-shell-ties-entropy-density-to-period}
The series has two ways of getting Einstein's equations for gravity. One borrows a standard action principle with a single strength of gravity. The other uses horizon thermodynamics: entropy proportional to area, a temperature set by acceleration, and heat flowing in as energy crosses a horizon. On the thermodynamic route alone, two ingredients were known to be independent assumptions: that the entropy per unit area does not depend on what matter is present, and the exact temperature of small horizons. This result asks what happens when both routes are kept, as the series now does. Wherever matter crosses a horizon, the heat balance then becomes a single equation tying the entropy per area to the temperature factor and the strength of gravity. Once the entropy per area is fixed at its standard value in the same units as the strength of gravity, either assumption gives the other there, and the known counterexamples stop working. The core equation is Jacobson's 1995 link between horizon entropy and the strength of gravity. In empty regions the equation says nothing, and the temperature remains a separate assumption there.
\end{theorem}
\noindent\emph{Rests on:} \ref{c:synthesis.macdowell-mansouri-principle-fixes-the-trace-free-part}, \ref{c:synthesis.imported-einstein-cartan-matter-coupling}, \ref{c:synthesis.r-loc-independent-sets-matter-coupling}, \ref{c:synthesis.einstein-equations-from-c-loc-and-clausius}, \ref{c:synthesis.integrated-loc-curvature-clause-forced-matter-clause-independent}, \ref{c:synthesis.loc-independent-of-clausius-and-killing-data}, \ref{c:synthesis.local-normality-at-the-edge-gives-the-null-ground-state}, \ref{c:synthesis.euler-entropy-of-a-horizon-patch-is-an-edge-term}, \ref{c:synthesis.imported-gtg-dictionary}, \ref{c:synthesis.imported-jacobson-clausius-1995}.

Both routes are stated for every gravitational field, so it matters which fields the series' quantum structure can reach. That turns on the tidal part of the curvature.

In four dimensions the tidal (Weyl) part of the curvature splits into two duality halves, the eigenspaces of multiplication by the unit pseudoscalar. For a real gravitational field the halves are conjugate, so if one vanishes so does the other, and the derived metric is conformally flat. Real black holes such as Schwarzschild and Kerr are therefore never half-flat.

\begin{proposition}[No real Lorentzian vacuum is half-flat]\label{c:correspondence.no-real-lorentzian-vacuum-half-flat}
In four dimensions the part of the gravitational field strength that describes tidal forces splits into two mirror-image halves. For a real gravitational field with one time direction the two halves are complex conjugates, so if one vanishes so does the other, and the geometry the field describes is then flat up to a change of scale. Real black holes such as Schwarzschild and Kerr are therefore never half-flat, which is why Penrose's construction of such spaces from spinor geometry (the nonlinear graviton) is stated only in complex or Euclidean settings.
\emph{First proved in} T9, Prop 6.5 (\href{https://doi.org/10.5281/zenodo.22722290}{doi:10.5281/zenodo.22722290}).
\end{proposition}
\noindent\emph{Rests on:} \ref{c:correspondence.imported-weyl-commutes-duality}, \ref{c:correspondence.imported-weyls-theorem}, \ref{c:correspondence.lorentzian-duality-complex-structure}, \ref{c:interior.exact-scale-decoupling}, \ref{c:interior.imported-kerr-curvature}.

The series attaches its quantum particles to spacetime through twistors, and twistor transport is path independent exactly where the Weyl curvature vanishes. On those fields Einstein's equations allow only perfect fluids of uniform rest density and constant-curvature vacuum, while Nordstr\"om's trace equation can be solved for any matter.

So consistency on the fields where the twistors attach does not single out Einstein's equations. Nordstr\"om's theory, ruled out by observation because it deflects no light, passes every consistency check there.

\begin{theorem}[On the gravitational fields where twistor transport is path independent, exactly those with vanishing Weyl curvature, Einstein's equations admit no perfect fluid whose density varies on its rest spaces and no vacuum field but constant curvature, while Nordstr\"om's trace equation is solvable for every stress trace with no condition on the matter; so the consistency conditions of a coupled system on these fields (Bianchi identities, conservation, second order, the Newtonian limit) do not give Einstein's equations without the clause that the whole stress map is the source]\label{c:synthesis.twistor-flat-fields-give-nordstrom-not-einstein}
The series attaches its quantum particles to spacetime through twistors, a space that is the same at every point only where the gravitational field has no tidal part of a particular kind (the Weyl curvature, the part that bends light around a star and travels as gravitational waves). Such fields are rescalings of flat spacetime. If Einstein's equations hold on them, the matter they contain must be spread evenly in its own rest frame: a star or a lumpy dust cloud cannot sit in such a field, and empty space can curve only uniformly, so the space outside a body is not of this kind. On the same fields an older theory of gravity, Nordstr\"om's, in which only the trace of the matter's stress-energy sources gravity, works for any matter at all, satisfies every consistency requirement, and reproduces Newton's law. It is ruled out by observation (it bends no light, and it gives the wrong advance of Mercury's perihelion, minus one sixth of Einstein's), so it is a counterexample to the logic, not a viable theory. So requiring the coupled system of matter and gravity to be consistent on the fields where the twistors attach does not single out Einstein's equations; something more is needed: the rule that the whole stress-energy, not only its trace, is the source, or an attachment of the particles that works where the Weyl curvature is nonzero.
\end{theorem}
\noindent\emph{Rests on:} \ref{c:synthesis.gtg-tractor-gauge-field-reproduces-tractor-connection}, \ref{c:correspondence.imported-tractor-connection-local-twistor-transport}, \ref{c:correspondence.imported-weyls-theorem}, \ref{c:synthesis.imported-gtg-dictionary}, \ref{c:synthesis.imported-lovelock-and-anderson-duchamp-natural-lagrangians}, \ref{c:synthesis.imported-nordstrom-scalar-theory-einstein-fokker-form}, \ref{c:field-theory.twistor-attachment}.

What Nordstr\"om's theory lacks is Einstein's rule that the whole stress, not only its trace, is the source. This step asks whether a symmetry supplies that rule.

At first order, at each non-null momentum, there are exactly four classes of gravitational fields closed under displacements. Varying only within the Weyl-flat class, the attachment's domain, couples gravity to the stress trace alone and decouples the Maxwell field, so gauge invariance does not give the rule. The missing part is the trace-free channel, exactly what the null fluxes $k \cdot T(k)$ of the Clausius route see. Under an action principle the rule amounts to varying the field off the attachment's domain.

\begin{theorem}[At first order, at each non-null momentum, the classes of gravitational fields that are closed under displacements and covariant under the stabiliser of the momentum are exactly four (at null momentum there are further classes, among them gauge plus the two helicity-two directions); varying the field only within the Weyl-flat class, the domain of the twistor attachment, couples it to the stress trace alone, and the Maxwell field not at all, so displacement gauge invariance does not give the whole-stress clause; the clause is the trace-free channel, which the null fluxes see exactly and which with conservation fixes a conserved source at every nonzero momentum]\label{c:synthesis.weyl-flat-variations-carry-only-the-stress-trace-and-the-whole-stress-clause-is-the-trace-free-channel}
Einstein's theory says that everything about matter's energy, momentum and pressure is a source of the gravitational field. An older rival, Nordstr\"om's, used only one number made from them, their trace, and is ruled out because it does not deflect light. The series wants to derive Einstein's rule from its own field theory. This result shows what that requires. If gravity's field is only allowed to change in the ways that keep the derived metric a rescaling of the flat one, which are exactly the fields where the series' twistors can be attached, then it responds only to the trace, as in Nordstr\"om's theory, and light and other massless fields do not affect it at all, even though this restricted theory respects every change of coordinates. The full rule needs the field to be able to change in five further ways that leave those fields. (For waves moving at the speed of light the bookkeeping has a few more cases, but the conclusion is the same.) The light-ray form of Einstein's equations used elsewhere in the series sees exactly the part that the trace misses, and energy conservation then fills in the rest, up to a cosmological constant. So the missing step is not a symmetry; it is that the matter theory must work, and be varied, beyond the fields where the twistors attach.
\end{theorem}
\noindent\emph{Rests on:} \ref{c:synthesis.twistor-flat-fields-give-nordstrom-not-einstein}, \ref{c:synthesis.imported-maxwell-stress-spacetime-algebra}, \ref{c:synthesis.imported-nordstrom-scalar-theory-einstein-fokker-form}, \ref{c:synthesis.imported-lovelock-and-anderson-duchamp-natural-lagrangians}, \ref{c:synthesis.imported-gtg-dictionary}.

The last question on this side is the strength of gravity itself.

In Sakharov's induced gravity the series' content produces $1/G$ from its quantum fluctuations; with a proper-time cut-off $\Lambda$, here a cut-off and not the cosmological constant, \[ 1/G = \bigl[B\Lambda\^{}{2} - \textstyle\sum w\, m\^{}{2}\ln(\Lambda\^{}{2}/m\^{}{2})\bigr]/(12\pi) \] If only the weak-scale content lies below the cut-off, the count $B$ vanishes and the remainder has the wrong sign and is far too small. Where $B > 0$ the observed $G$ can be matched, but only by choosing the cut-off to match it, which is a fit. On this route $G$ stays an input.

\begin{theorem}[In the Feynman gauge with a proper-time cut-off $\Lambda$ the series' content induces $1/G = [B\Lambda^{2} - \sum w m^{2}\ln(\Lambda^{2}/m^{2})]/(12\pi) + O(m^{2})$, with $B = -36$ plus the scalars' weights for the chain's Pati-Salam content under Hypothesis (S); the logarithmic remainder is repulsive if the content below $\Lambda$ is the weak-scale content alone with $B = 0$, its sign is not fixed by the store for a cut-off above $M_{I}$ under (S), and it stays far below $1/G$; where $B > 0$ matching $G$ fixes $\Lambda$, so the route gives $G$ only by choosing the cut-off]\label{c:synthesis.series-content-induces-a-newton-coupling-only-through-an-unfixed-cut-off}
Sakharov's idea is that gravity's strength could be produced by the quantum jitter of all other fields, up to some shortest length where the jitter is cut off. The result is the cut-off energy squared times a count of the fields, plus a much smaller term from the particle masses. If the series' larger unified symmetry is real at high energies, the count depends on how its many extra Higgs fields couple to curvature, which nothing in the series decides. If only the everyday particles lie below the cut-off and their count is zero, the small mass term is left, and it has the wrong sign: it would make gravity repulsive, and it is far too weak. With the series' unified symmetry, heavy force carriers pull the other way and the sign is undecided, but the term is still far too weak. Where the count is positive, gravity's observed strength can be matched, but only by choosing the cut-off to match it, a little above the Planck energy. The series has no number that would set that cut-off, so on this route the strength of gravity remains an input.
\end{theorem}
\noindent\emph{Rests on:} \ref{c:synthesis.imported-sakharov-induced-gravity-coefficients}, \ref{c:synthesis.attachment-modules-give-sakharov-weights-cancelling-at-three-separate-generations}, \ref{c:standard-model.connection-forced}, \ref{c:standard-model.chain}, \ref{c:synthesis.linear-rung-field-has-a-neutral-normal-pair-in-the-scalon-sum}, \ref{c:synthesis.krein-quartet-drops-from-the-induced-einstein-hilbert-coefficient-only-under-a-regulator-that-treats-it-alike}.

\section{Temperature and Hawking radiation}

This part follows the temperature that the horizon route needs back to a statement about the quantum state, and asks how much of that statement is now proved.

The temperature clause can be traded for a premise of a different kind. Premise (N$_{+}$) says only that the modes on a horizon's light rays are in a ground state for sliding along the rays, with no temperature in it.

Seen from one side of the edge, that ground state is exactly thermal at rapidity period $2\pi$: the period comes out of positive energy and of the observer seeing only half of each ray. So the continuum forms of (R) for the eternal hole, of (R$_{\Lambda}$) and of (R$_{\mathrm{loc}}$) in Minkowski space all follow from one premise. The input that remains is regularity across the bifurcation surface, where the horizon's two halves meet.

\begin{theorem}[A ground state for the affine null translations of a horizon's generators restricts on either side of the edge to the continuum I-state at rapidity period $2\pi$, so the continuum restatements of (R) for the eternal hole, (R$_{\Lambda}$), (R$_{\mathrm{loc}}$) and the joint KMS part of (Q) follow from one premise with no temperature in it, posable only on the continuum]\label{c:synthesis.positive-null-energy-gives-r-premises}
Black holes, the horizon of an expanding universe, and the small horizons an accelerating observer sees all have a temperature. The series has so far assumed each of these temperatures with a separate premise of the same kind, each fixing a period of $2\pi$ by hand. This result shows that a single assumption with no temperature in it gives all three, once the quantum modes are taken to run continuously along each horizon's light rays. The assumption is that these modes are in their lowest-energy state with respect to sliding along the rays. Looked at from one side of the horizon, that state is exactly thermal, and the $2\pi$ comes out of the calculation, because energy is positive and the observer sees only half of each ray. There are costs. The finite lists of modes the series counts become an approximation. For a black hole the assumption describes an eternal hole in equilibrium, not one formed by collapse. And the assumption is still an input: it says the state is smooth across the point where the horizon's two halves meet, which is what every derivation of black-hole temperature assumes in some form.
\end{theorem}
\noindent\emph{Rests on:} \ref{c:hawking.i-state-continuum}, \ref{c:hawking.imported-quasi-free-kms-reduces-two}, \ref{c:hawking.imported-ideal-fermi-gas-uniqueness}, \ref{c:singularity.euclidean-section-period}, \ref{c:synthesis.de-sitter-static-boost-frozen-point}, \ref{c:synthesis.local-boost-temperature-is-local-regularity}, \ref{c:synthesis.point-spectrum-boost-admits-no-half-sided-translations}, \ref{c:angle.imported-theorems-deliver-their-hypotheses}, \ref{c:mass.chart-form-spin-isomorphism}, \ref{c:synthesis.n-plus-null-ground-state}.

Where gravity is present, what the horizon route uses is a local form of (N$_{+}$), at each point of each local horizon.

On a single generator, the light-ray version of the Hadamard condition holds exactly when the state differs from the null vacuum by a smooth kernel. With no stationarity and no temperature, that gives the local form of (N$_{+}$) at every point. It does not single out the vacuum, since thermal states pass it too.

\begin{theorem}[On a horizon generator the generator analogue of the microlocal (Hadamard) condition, on both orderings of the two-point kernel, is equivalent to a smooth difference from the null vacuum, so with no stationarity and no temperature it gives the local form of (N$_{+}$) at every point (with the stored theorem and boost stationarity at the edge, (N$_{+}$)); the affine-KMS states satisfy it, so it does not select the ground state]\label{c:synthesis.microlocal-condition-on-a-generator-gives-the-local-null-ground-state}
Physicists usually require that any reasonable quantum state, looked at closely enough, resembles empty space. The precise version of this requirement says that the very fine-grained structure of the state carries only positive energy. This result writes a version of that requirement for the quantum modes running along a single light ray of a horizon and shows that it is the same as saying the state differs from the vacuum only by something smooth. From there the store's earlier results give the horizon temperature at every point, without assuming that anything is static. It does not single out the vacuum, though: a gas of modes at any temperature also passes, so to get the full premise one still needs the state to be static near the point where the two halves of the horizon meet. Whether the series' states obey the requirement is still an open assumption.
\end{theorem}
\noindent\emph{Rests on:} \ref{c:synthesis.local-normality-at-the-edge-gives-the-null-ground-state}, \ref{c:synthesis.stationary-horizon-state-is-its-own-edge-scaling-limit}, \ref{c:synthesis.positive-null-energy-gives-r-premises}, \ref{c:synthesis.imported-microlocal-hadamard-condition}, \ref{c:synthesis.imported-local-quasi-equivalence-of-hadamard-states}, \ref{c:synthesis.imported-kay-wald-uniqueness-and-no-go}, \ref{c:synthesis.imported-propagation-of-hadamard-form}, \ref{c:synthesis.imported-powers-stormer-takesaki-disjointness}.

The question becomes whether physically reasonable states satisfy that condition on the light rays. Where there is no gravitational field, the answer is yes.

On the flat field every gauge-invariant quasi-free Hadamard state of the Dirac field restricts on every null plane to the generator condition. So the local form of (N$_{+}$) holds at every point of every local Rindler horizon there.

\begin{theorem}[On the flat gravitational field every gauge-invariant quasi-free Dirac state differing from a vacuum by a smooth kernel restricts on every null plane to a state satisfying the generator microlocal condition, the vacuum itself to the null vacuum, so it has the local form of (N$_{+}$) on every local Rindler horizon; not the ground-state clause, not nonzero field strength]\label{c:synthesis.bulk-hadamard-states-restrict-to-the-generator-condition-on-flat-null-planes}
The series' account of horizon temperature rests on an assumption that, very close to any point of a horizon, the quantum field looks like empty space. This result checks that assumption against the field theory in empty flat space. On any light-like plane, the field's vacuum, whatever the particle's mass, looks exactly like the empty state that the assumption names, and every state that physicists regard as reasonable, one that differs from the vacuum only smoothly, looks like it at short distances too. So in flat space the assumption's short-distance part is a consequence of the field theory, though not the part that says the state is the vacuum everywhere, and not yet in a gravitational field.
\end{theorem}
\noindent\emph{Rests on:} \ref{c:synthesis.light-ray-good-part-of-a-spinor-is-two-chirality-arrays}, \ref{c:synthesis.slot-module-light-ray-arrays-are-weyl-components}, \ref{c:synthesis.positive-null-energy-gives-r-premises}, \ref{c:field-theory.second-quantisation}, \ref{c:synthesis.microlocal-condition-on-a-generator-gives-the-local-null-ground-state}, \ref{c:synthesis.imported-microlocal-hadamard-condition}, \ref{c:synthesis.imported-null-plane-restriction-of-free-fields-is-mass-independent}.

With gravity present the work splits into two parts, the bulk and the light rays.

In the bulk, a gauge-invariant quasi-free Hadamard state of the Dirac field of any mass, looked at closely enough at any point of a globally hyperbolic gravitational field, becomes the flat massless vacuum. This is a standard theorem of Sahlmann and Verch, redone in the spacetime algebra: the mass, the connection and the curvature fade out of the correlations at definite rates.

\begin{theorem}[Given the Hadamard parametrix of the squared Dirac operator (imported), every gauge-invariant quasi-free Hadamard state of the Dirac field of mass m0 >= 0 on a globally hyperbolic gravitational field scales, under the store's dilation in a gauge with hb(x0) = 1, to the flat field's massless vacuum on both orderings (standard: Sahlmann and Verch 2001, Prop. 5.6, re-derived in the spacetime algebra for Dirac spinors), so clause (a) of (HS) holds and the premise's content is its generator clause (b)]\label{c:synthesis.hadamard-dirac-states-scale-to-the-flat-massless-vacuum-at-every-point-of-a-gravitational-field}
The series assumed that any state of the free electron field of the standard regular kind (gauge-invariant, quasi-free and Hadamard), looked at closely enough at a point where gravity is present, resembles the empty state of flat space. This result proves that assumption for the bulk of space-time, for massive and massless particles and for every gravitational field with a well-posed causal structure (globally hyperbolic), by redoing in the series' algebra a known theorem of Sahlmann and Verch from the mathematical physics literature: close up, the curvature, the mass and the gravitational connection all fade out of the state's correlations at definite rates. It relies on two standard analytic facts quoted from the literature rather than proved in the series. What it does not prove is the companion statement the series also needs: that the same is true when the state is looked at only along the light rays of a small horizon. That part stays an assumption.
\end{theorem}
\noindent\emph{Rests on:} \ref{c:synthesis.imported-gtg-dictionary}, \ref{c:synthesis.imported-dimock-dirac-car-algebra-on-a-gravitational-field}, \ref{c:field-theory.second-quantisation}, \ref{c:synthesis.imported-microlocal-hadamard-condition}, \ref{c:synthesis.imported-dirac-hadamard-form-propagates-and-adiabatic-states-are-hadamard}, \ref{c:synthesis.imported-hadamard-parametrix-of-the-squared-dirac-operator-leads-with-van-vleck-root-times-spin-transport}, \ref{c:synthesis.imported-free-field-spin-statistics-and-guido-longo-theorem}, \ref{c:synthesis.bulk-hadamard-states-restrict-to-the-generator-condition-on-flat-null-planes}.

On the light rays the danger is that gravity's tidal stretching of the rays leaves a mark at the shortest distances.

At leading order it does not: on a null hypersurface the Weyl and Ricci field strength along the generators, the expansion and the shear cancel exactly from the leading singularity, because \[ \Delta\^{}{1/2}/\sqrt{\det M} = (a(u)\,a(u'))\^{}{-1/2} \] Hadamard states differ there by smooth kernels, so the question no longer depends on the state, only on the geometry.

\begin{theorem}[On every null hypersurface of a gravitational field generated by affine null geodesics, the field strength along the generators, the expansion and the shear cancel exactly from the coefficient of the leading null singularity of a Hadamard Dirac kernel restricted to it, which at leading order in the canonical normalisation is the null vacuum; restricted kernels of Hadamard states differ by smooth kernels, so clause (b) of (HS) is a property of the field and reduces to a remainder bound on the Hadamard parametrix, argued but not proved]\label{c:synthesis.field-strength-cancels-from-the-leading-null-singularity-of-hadamard-dirac-states-on-null-hypersurfaces}
The series assumes that, looked at closely along the light rays that make up a small horizon, any reasonable quantum state of the electron field resembles empty space, even where gravity is present. One way this could fail is if gravity's tidal stretching of the light rays left a mark on the state at the shortest distances. This result shows that, at the leading order, it does not: the tidal effects change how light rays spread and how distances between them are measured, but the two changes cancel exactly. It also shows that the assumption is either true for all reasonable states or for none, so it is really a question about gravity's geometry, not about which state the world is in. Smaller correction terms needed a technical estimate, argued here and proved by a later result.
\end{theorem}
\noindent\emph{Rests on:} \ref{c:synthesis.imported-gtg-dictionary}, \ref{c:synthesis.imported-microlocal-hadamard-condition}, \ref{c:synthesis.imported-dirac-hadamard-form-propagates-and-adiabatic-states-are-hadamard}, \ref{c:synthesis.imported-hadamard-parametrix-of-the-squared-dirac-operator-leads-with-van-vleck-root-times-spin-transport}, \ref{c:synthesis.imported-dimock-dirac-car-algebra-on-a-gravitational-field}, \ref{c:synthesis.bulk-hadamard-states-restrict-to-the-generator-condition-on-flat-null-planes}, \ref{c:synthesis.microlocal-condition-on-a-generator-gives-the-local-null-ground-state}.

Two smaller corrections could each have grown without limit.

They cancel: on every local Rindler horizon, for every mass, the restricted Dirac Hadamard parametrix is the null vacuum $W_{0}\otimes\id_{K}$ plus a bounded kernel, the $\log\epsilon$ terms cancelling identically. So the remainder has generator scaling degree at most $0$.

\begin{theorem}[On every piece of a geodesic null hypersurface of a gravitational field whose distinct generators are spacelike separated (every small enough piece, every local Rindler horizon), the Dirac Hadamard parametrix, restricted to the generators in the canonical normalisation as the epsilon-regularised transverse-smeared limit, is the null vacuum plus a bounded kernel: the log epsilon terms cancel identically, so the remainder bound posed by (HS) holds for every mass]\label{c:synthesis.restricted-dirac-hadamard-parametrix-on-a-null-hypersurface-is-the-null-vacuum-plus-a-bounded-kernel}
The series needs every reasonable quantum state of the electron field to look like empty space along the light rays of the small horizons used to derive the equations of gravity, when looked at closely enough. The question had been reduced to a statement about gravity's geometry alone: a fixed mathematical object, which carries the short-distance behaviour shared by all reasonable states, should look like that of empty space along those light rays, up to corrections that fade at short distances. The strongest part had been checked before; two weaker pieces, which separately would grow without limit, were expected but not shown to cancel. This result shows that they do cancel exactly, on small enough pieces of the light sheets of every gravitational field (including the small horizons the series uses) and for every mass of the field, and that everything left over stays bounded. So the assumed statement is now proved, for the precise way the object is restricted to the light rays. One step remains before it applies to the states themselves: showing that this way of restricting agrees with the more fundamental one.
\end{theorem}
\noindent\emph{Rests on:} \ref{c:synthesis.imported-gtg-dictionary}, \ref{c:synthesis.imported-hadamard-parametrix-of-the-squared-dirac-operator-leads-with-van-vleck-root-times-spin-transport}, \ref{c:synthesis.imported-dirac-hadamard-form-propagates-and-adiabatic-states-are-hadamard}, \ref{c:synthesis.field-strength-cancels-from-the-leading-null-singularity-of-hadamard-dirac-states-on-null-hypersurfaces}, \ref{c:synthesis.bulk-hadamard-states-restrict-to-the-generator-condition-on-flat-null-planes}.

That computation used one particular way of restricting to the light rays, so the restriction itself has to be put on firm ground.

The Dirac current is future causal, so the flux through a null sheet never exceeds the total: $W^{+} + W^{-} = \iota_{\mathcal{H}}^{*}\iota_{\mathcal{H}} \le 1$. Hence every state of the Dirac field restricts to the generators with no characteristic Cauchy problem, and Hadamard states differ there by smooth kernels. If the anticommutation relations hold on the generators, the restriction is the limit of the bulk two-point function on thinner and thinner layers.

\begin{theorem}[Every state of the Dirac field on a gravitational field restricts to a null hypersurface through the flux contraction of the conserved current, with no characteristic Cauchy problem, its restricted kernels have generator scaling degree at most one, and restricted kernels of Hadamard states differ by smooth pullbacks; if the generator anticommutation relations hold on nearby null hypersurfaces (proved in the literature for light cones), the restriction is the thin-layer limit of the bulk two-point function, and the restricted Hadamard parametrix then has a regularisation-free definition, whose agreement with the regularised one used in (HS) is argued]\label{c:synthesis.dirac-states-restrict-to-null-hypersurfaces-through-the-flux-contraction-and-are-thin-layer-limits-under-the-generator-car}
The series assumes that, seen closely along the light rays of a small horizon, quantum states of the electron field look like empty space. Testing this needs a precise way to look at a state "on the light rays only", and the earlier work borrowed one from memory. This result builds it from a basic fact: the flow of electron charge-current is never faster than light, so what flows through a sheet of light rays can never exceed the total. That alone makes every state restrict to the light rays, caps how singular the restriction can be, and shows that any two reasonable states differ there only by a smooth amount. If the light-ray field also obeys the usual rules for fermions, which mathematicians have proved for light cones but which are only believed for general light sheets, then the restriction is simply the limit of ordinary measurements made in thinner and thinner slabs, and the earlier borrowed step is no longer needed. Whether this matches the way the horizon assumption itself is computed is argued, not proved.
\end{theorem}
\noindent\emph{Rests on:} \ref{c:synthesis.imported-gtg-dictionary}, \ref{c:synthesis.imported-dimock-dirac-car-algebra-on-a-gravitational-field}, \ref{c:synthesis.imported-microlocal-hadamard-condition}, \ref{c:synthesis.imported-dirac-hadamard-form-propagates-and-adiabatic-states-are-hadamard}, \ref{c:synthesis.imported-hadamard-parametrix-of-the-squared-dirac-operator-leads-with-van-vleck-root-times-spin-transport}, \ref{c:synthesis.imported-haefner-nicolas-dirac-characteristic-problem-on-light-cones-near-the-vertex}, \ref{c:synthesis.field-strength-cancels-from-the-leading-null-singularity-of-hadamard-dirac-states-on-null-hypersurfaces}, \ref{c:synthesis.bulk-hadamard-states-restrict-to-the-generator-condition-on-flat-null-planes}, \ref{c:synthesis.microlocal-condition-on-a-generator-gives-the-local-null-ground-state}.

Those anticommutation relations were the remaining condition.

On every geodesic null hypersurface the Dirac causal propagator between the sheet and a nearby leaf tends to the canonical delta kernel on the generators, the one possible correction vanishing by the same half-density law. So the generator anticommutation relations hold, local Rindler horizons included.

\begin{theorem}[On every geodesic null hypersurface of a gravitational field the causal propagator of the Dirac field, restricted between the hypersurface and a nearby leaf of a null foliation, tends to the canonical delta kernel of the generators, the step term cancelling by the half-density law; with the flux contraction of the conserved current, the generator anticommutation relations hold there, and the restricted Hadamard parametrix in the thin-layer definition is the null vacuum plus a bounded kernel]\label{c:synthesis.dirac-causal-propagator-between-nearby-null-leaves-tends-to-the-generator-car-on-every-geodesic-null-hypersurface}
Earlier work needed a basic rule for the electron field seen only along the light rays of a small horizon: the field there should obey the standard rules for fermions, exactly as in empty space. Mathematicians had proved this for light cones; for the general light sheets used in the series it was only believed. This result proves it for every such sheet by a short calculation with the one object every state shares, the propagator that carries signals at or below light speed. Two corrections that could have spoiled the rule turn out to cancel, because the area of the sheet and the focusing of light rays are tied together. Together with two companion results, it also shows that every reasonable state looks like empty space along these light rays at short distances.
\end{theorem}
\noindent\emph{Rests on:} \ref{c:synthesis.imported-gtg-dictionary}, \ref{c:synthesis.imported-dimock-dirac-car-algebra-on-a-gravitational-field}, \ref{c:synthesis.imported-hadamard-parametrix-of-the-squared-dirac-operator-leads-with-van-vleck-root-times-spin-transport}, \ref{c:synthesis.imported-dirac-hadamard-form-propagates-and-adiabatic-states-are-hadamard}, \ref{c:synthesis.imported-microlocal-hadamard-condition}, \ref{c:synthesis.field-strength-cancels-from-the-leading-null-singularity-of-hadamard-dirac-states-on-null-hypersurfaces}, \ref{c:synthesis.bulk-hadamard-states-restrict-to-the-generator-condition-on-flat-null-planes}, \ref{c:synthesis.dirac-states-restrict-to-null-hypersurfaces-through-the-flux-contraction-and-are-thin-layer-limits-under-the-generator-car}, \ref{c:synthesis.restricted-dirac-hadamard-parametrix-on-a-null-hypersurface-is-the-null-vacuum-plus-a-bounded-kernel}.

The pieces now assemble into one theorem.

(HS), formerly a premise, is proved: for every gauge-invariant quasi-free Hadamard state of the Dirac field of any mass, at every point of every local Rindler horizon of every gravitational field, the scaling limit on the generators is the null vacuum. That is the local form of (N$_{+}$) at every such point. It rests on imported standard results and on a few standard inputs cited from memory, such as a local null foliation and the Morse lemma with parameters, and it is stated for the free Dirac field only. It does not give the ground-state clause of (N$_{+}$), nor (Loc), nor which state the world is in.

\begin{theorem}[(HS), the Hadamard scaling limit, a theorem (formerly a premise): for every gauge-invariant quasi-free Hadamard state of the Dirac field of any mass on a gravitational field, at every point of every local Rindler horizon, the bulk scaling limit is the flat massless vacuum and the restriction to the generators through the flux trace has the null vacuum as its generator scaling limit, the epsilon-regularised restriction of the parametrix agreeing with it modulo a bounded time-function-dependent kernel; hence the local form of (N$_{+}$) at every such point for these states]\label{c:synthesis.hadamard-scaling-limit-premise}
Physicists require that any reasonable quantum state, looked at closely enough, resembles empty space. The series had assumed this along the light rays of the small horizons used to derive the equations of gravity, wherever gravity is present. It is now proved, for the electron field of any mass, given standard results from the literature: the flow of electric charge through a sheet of light can never exceed the total, signals travel at or below light speed, and two growing corrections cancel because the area of the sheet and the focusing of its light rays are tied together. So every reasonable state looks like empty space along those light rays at short distances, which is what the series' horizon temperature needs locally. It does not pick out which state the world is in, and it does not by itself give the full horizon assumption, since warm states pass the test too.
\end{theorem}
\noindent\emph{Rests on:} \ref{c:synthesis.imported-dimock-dirac-car-algebra-on-a-gravitational-field}, \ref{c:synthesis.imported-microlocal-hadamard-condition}, \ref{c:synthesis.imported-dirac-hadamard-form-propagates-and-adiabatic-states-are-hadamard}, \ref{c:synthesis.imported-hadamard-parametrix-of-the-squared-dirac-operator-leads-with-van-vleck-root-times-spin-transport}, \ref{c:synthesis.hadamard-dirac-states-scale-to-the-flat-massless-vacuum-at-every-point-of-a-gravitational-field}, \ref{c:synthesis.field-strength-cancels-from-the-leading-null-singularity-of-hadamard-dirac-states-on-null-hypersurfaces}, \ref{c:synthesis.restricted-dirac-hadamard-parametrix-on-a-null-hypersurface-is-the-null-vacuum-plus-a-bounded-kernel}, \ref{c:synthesis.dirac-states-restrict-to-null-hypersurfaces-through-the-flux-contraction-and-are-thin-layer-limits-under-the-generator-car}, \ref{c:synthesis.dirac-causal-propagator-between-nearby-null-leaves-tends-to-the-generator-car-on-every-geodesic-null-hypersurface}, \ref{c:synthesis.bulk-hadamard-states-restrict-to-the-generator-condition-on-flat-null-planes}.

The ground-state clause is the part of the null ground-state premise that this leaves open.

Asking the state to look like the null vacuum near every point of a generator, including the point at infinity where the vacuum's symmetries close the generator up, puts it in the null vacuum's family, its folium, with positive translations along the ray. Being the null vacuum itself, which is (N$_{+}$), still needs translation or boost stationarity. Normality at finite points alone does not even fix the family: a state with a kink at infinity passes it.

\begin{theorem}[On a horizon generator closed by its Moebius covariance, the Hilbert-Schmidt form of local normality to the null vacuum on every interval, the arcs through the point at infinity included, is a globally square-integrable two-point difference; for pure gauge-invariant quasi-free states this is global quasi-equivalence with the null vacuum, the folium form of (N$_{+}$) with the affine translations positive in the state's representation, and it gives (N$_{+}$) itself only with translation or boost stationarity; it excludes the affine-KMS states on one far end, while normality at every finite point, the vacuum on each far end separately and implementability of the affine group do not fix the folium]\label{c:synthesis.normality-on-the-closed-generator-gives-the-vacuum-folium-not-the-ground-state}
The series assumes that the quantum field on the light rays of a horizon is in its lowest-energy state. Earlier work showed that for a field that does not change in time this is the same as a standard locality requirement at the point where the horizon's light rays begin. Here the time-independence is dropped. Looking like the vacuum near every point of a light ray is not enough, since a warm state does so too; asking for the vacuum far along the ray rules out the warm state, but a subtler state still passes if the two far ends are tested separately. Joining the two ends of a light ray at one point, as the symmetries of the vacuum allow, and requiring the vacuum's local behaviour there as well puts the state in the vacuum's family, where energy along the ray is never negative; being the vacuum itself still needs time-independence.
\end{theorem}
\noindent\emph{Rests on:} \ref{c:synthesis.local-normality-at-the-edge-gives-the-null-ground-state}, \ref{c:synthesis.imported-powers-stormer-takesaki-disjointness}, \ref{c:synthesis.stationary-horizon-state-is-its-own-edge-scaling-limit}, \ref{c:synthesis.positive-null-energy-gives-r-premises}, \ref{c:synthesis.attachment-conformal-energy-is-generated-by-z}, \ref{c:synthesis.collapsing-hole-flux-from-ground-state-on-the-ingoing-sheet}, \ref{c:synthesis.imported-haag-local-normality}.

\section{Quantum fields in one algebra}

The last part asks where the series' quantum state could come from.

The natural candidate is the series' own attachment of quantum modes to spacetime through conformal spinors, which in flat space comes with a built-in positive energy.

On a gravitational field those spinors can be parallel for the tractor gauge field only where the Weyl curvature vanishes: a two-dimensional subspace at type N points, none at types I, II, D and III. No conformal-energy generator commutes with a nonzero Weyl curvature, so off Weyl-flat regions the attachment supplies no state. Real black holes, which are never half-flat, lie outside its reach.

\begin{theorem}[On a gravitational field the conformal spinors of the attachment are parallel for the tractor gauge field only where the Weyl curvature vanishes: at a type N point a two-dimensional subspace can be, at points of type I, II, D or III none, and no conformal-energy generator commutes with a nonzero Weyl curvature, so off Weyl-flat regions the tractor transport of the attachment supplies no state]\label{c:synthesis.attachment-transport-is-flat-only-where-weyl-vanishes}
The series links its quantum modes to spacetime through a special four-dimensional space of spinors, which in flat space comes with a built-in notion of positive energy. In a gravitational field one has to carry that space from point to point, and gravity's tidal part, the Weyl curvature, twists it on the way. This result shows that the space survives only where that tidal part vanishes: at a point of one special radiation type half of it can survive, and at every other point none of it, and the built-in notion of positive energy is lost. So in a general gravitational field this way of attaching the modes does not fix the quantum state; the state has to be supplied in some other way.
\end{theorem}
\noindent\emph{Rests on:} \ref{c:synthesis.gtg-tractor-gauge-field-reproduces-tractor-connection}, \ref{c:synthesis.imported-gtg-dictionary}, \ref{c:field-theory.twistor-attachment}, \ref{c:synthesis.attachment-conformal-energy-is-generated-by-z}, \ref{c:field-theory.where-each-lives}, \ref{c:synthesis.imported-twistor-equation-integrability-petrov-n-or-o}.

\section{Results used from elsewhere}
The proofs above use the following results of the series, stated briefly; each is proved in the paper named.

\begin{fact}[Imported: the theorems that deliver $2\pi$, and their hypotheses]\label{c:angle.imported-theorems-deliver-their-hypotheses}
Under the spectrum condition, covariance and locality, the modular flow of a wedge or half-sided algebra is the boost at rate $2\pi$, and a regular stationary quasi-free state at a bifurcate Killing horizon is KMS at $\kappa/2\pi$.
\emph{Source:} J. J. Bisognano and E. H. Wichmann, On the duality condition for a Hermitian scalar field, J. Math. Phys. 16 (1975), 985--1007.; J. J. Bisognano and E. H. Wichmann, On the duality condition for quantum fields, J. Math. Phys. 17 (1976), 303--321.; H.-J. Borchers, The CPT-theorem in two-dimensional theories of local observables, Commun. Math. Phys. 143 (1992), 315--332.; P. D. Hislop and R. Longo, Modular structure of the local algebras associated with the free massless scalar field theory, Commun. Math. Phys. 84 (1982), 71--85.; R. Brunetti, D. Guido and R. Longo, Modular structure and duality in conformal quantum field theory, Commun. Math. Phys. 156 (1993), 201--219.; G. L. Sewell, Quantum fields on manifolds: PCT and gravitationally induced thermal states, Ann. Phys. 141 (1982), 201--224.; B. S. Kay and R. M. Wald, Theorems on the uniqueness and thermal properties of stationary, nonsingular, quasifree states on spacetimes with a bifurcate Killing horizon, Phys. Rep. 207 (1991), 49--136.; W. G. Unruh, Notes on black-hole evaporation, Phys. Rev. D 14 (1976), 870--892..
\end{fact}

\begin{fact}[Imported: the tractor connection, and the conformal spinor connection (local twistor transport)]\label{c:correspondence.imported-tractor-connection-local-twistor-transport}
The conformal class of a gravitational field on flat spacetime carries the tractor connection, a $\Cl(2,4)$ gauge field, and its spinor form, the conformal spinor connection on the rank-four bundle of conformal spinors, with explicit formulas imported in the series' conventions.
\emph{Source:} T9, Fact 4.9 (\href{https://doi.org/10.5281/zenodo.22722290}{doi:10.5281/zenodo.22722290}).
\end{fact}

\begin{fact}[Imported: Weyl commutes with duality]\label{c:correspondence.imported-weyl-commutes-duality}
The Weyl map of the gravitational field strength commutes with duality, $\mathcal{W}(IB) = I\,\mathcal{W}(B)$: on an oriented Lorentzian four-manifold the Weyl operator on two-forms commutes with the Hodge star.
\emph{Source:} R. Penrose and W. Rindler, Spinors and Space-Time, Vols. 1--2, Cambridge University Press (1984, 1986)..
\end{fact}

\begin{fact}[Imported: Weyl's theorem]\label{c:correspondence.imported-weyls-theorem}
In dimension at least four a gravitational field on flat spacetime is locally a pure rescaling, $\hb^{-1}(a) = e^{\omega}a$ in some gauge, if and only if its Weyl curvature vanishes (Weyl 1918).
\emph{Source:} H. Weyl, Reine Infinitesimalgeometrie, Math. Z. 2 (1918), 384--411..
\end{fact}

\begin{fact}[Lorentzian duality is a complex structure]\label{c:correspondence.lorentzian-duality-complex-structure}
On $2$-forms of an oriented four-dimensional Lorentzian space the Hodge star squares to $-1$, with three-dimensional, mutually conjugate eigenspaces for the eigenvalues $\pm i$.
\emph{Source:} T9, Lem 6.4 (\href{https://doi.org/10.5281/zenodo.22722290}{doi:10.5281/zenodo.22722290}).
\end{fact}

\begin{fact}[Imported: the frozen point and its geometry (external ingredients)]\label{c:dark-energy.imported-frozen-point-geometry-external}
The frozen-point fact takes its tractor connection and almost-Einstein scale tractor from Gover's work on almost-Einstein manifolds.
\emph{Source:} A. R. Gover, Almost Einstein and Poincar'e-Einstein manifolds in Riemannian signature, J. Geom. Phys. 60 (2010), 182--204..
\end{fact}

\begin{fact}[Second quantisation]\label{c:field-theory.second-quantisation}
With the single-algebra field and $i = I_4$, second quantisation gives a local field and the standard $d\Gamma(M)$, one chirality entering complex-linearly and the other conjugated.
\emph{Source:} FT, Prop 5.2.
\end{fact}

\begin{fact}[The attachment through conformal spinors]\label{c:field-theory.twistor-attachment}
The modes attach to spacetime through the conformal spinor space, the palatial storey giving a positive-energy one-particle space and the Penrose transform the mode functions, with floors as energy cut-offs and local CAR lifts exactly for odd $2s$.
\emph{Source:} FT, Prop 3.2.
\end{fact}

\begin{fact}[Where each imaginary unit lives]\label{c:field-theory.where-each-lives}
In real geometric algebra the conformal spinors' unit is $\pm I_6 = \pm I_4E$, equal to $\pm I_4$ on the two chiral doublets, and a massless field's unit is $\pm I_4$, with which a mass term anticommutes.
\emph{Source:} FT, Prop 4.1.
\end{fact}

\begin{fact}[The I-state on the continuum]\label{c:hawking.i-state-continuum}
On the continuum mode space at $\beta = 2\pi/\kappa$, for any self-adjoint $h$, the quasi-free state with symbol $(1 + e^{\beta h})^{-1}$ is $\tau$-KMS and is the unique $\tau$-KMS state.
\emph{Source:} T7, Thm 4.2 (\href{https://doi.org/10.5281/zenodo.22722121}{doi:10.5281/zenodo.22722121}).
\end{fact}

\begin{fact}[Imported: ideal-Fermi-gas uniqueness]\label{c:hawking.imported-ideal-fermi-gas-uniqueness}
For free fermions with quasi-free dynamics the $\beta$-KMS state is unique and is the quasi-free state with Fermi-Dirac symbol $(1 + e^{\beta h})^{-1}$.
\emph{Source:} O. Bratteli and D. W. Robinson, Operator Algebras and Quantum Statistical Mechanics 2, 2nd ed., Springer, Berlin (1997)., Ex.~5.3.2 and \S 5.3; H. Araki, On quasifree states of CAR and Bogoliubov automorphisms, Publ. RIMS Kyoto Univ. 6 (1971), 385--442..
\end{fact}

\begin{fact}[Imported: quasi-free KMS reduces to the two-point level]\label{c:hawking.imported-quasi-free-kms-reduces-two}
A gauge-invariant quasi-free CAR state is KMS for quasi-free dynamics exactly when its two-point function satisfies the two-point KMS identity.
\emph{Source:} H. Araki, On quasifree states of CAR and Bogoliubov automorphisms, Publ. RIMS Kyoto Univ. 6 (1971), 385--442.; O. Bratteli and D. W. Robinson, Operator Algebras and Quantum Statistical Mechanics 2, 2nd ed., Springer, Berlin (1997)., \S 5.3.
\end{fact}

\begin{fact}[Exact scale decoupling]\label{c:interior.exact-scale-decoupling}
The Schwarzschild curvature operator is exactly $-(M/r^3)\,S$ with a fixed, trace-free shape operator $S$ satisfying $S^2 + S = 2$, so only the scale changes on the approach to the singularity.
\emph{Source:} T6, Prop 3.3 (\href{https://doi.org/10.5281/zenodo.22721990}{doi:10.5281/zenodo.22721990}).
\end{fact}

\begin{fact}[Imported: the Kerr curvature]\label{c:interior.imported-kerr-curvature}
The Kerr field strength on physical bivectors is $\mathcal{R}(B) = \tfrac{M}{2}(r - Ia\cos\theta)^{-3}(B + 3\sigma_r B\sigma_r)$, with the inverse cube taken in $\R \oplus \R I$; imported from the literature.
\emph{Source:} C. Doran and A. Lasenby, Geometric Algebra for Physicists, Cambridge University Press (2003)., Ch.~14; B. Carter, Global structure of the Kerr family of gravitational fields, Phys. Rev. 174 (1968), 1559--1571..
\end{fact}

\begin{fact}[Imported: the tractor connection]\label{c:interior.imported-tractor-connection}
For a gravitational field on flat spacetime, the tractor bundle of its conformal class (the $\Cl(2,4)$ vectors $\sigma\nInf + \mu + \rho\nO$) carries a canonical metric-preserving connection, given in a chosen scale by an explicit formula involving the Schouten tensor.
\emph{Source:} T. N. Bailey, M. G. Eastwood and A. R. Gover, Thomas's structure bundle for conformal, projective and related structures, Rocky Mountain J. Math. 24 (1994), 1191--1217.; S. N. Curry and A. R. Gover, An introduction to conformal geometry and tractor calculus, with a view towards analysis on manifolds, in Asymptotic Analysis in General Relativity, London Mathematical Society Lecture Note Series 443, Cambridge University Press (2018), 86--170., \S 5.
\end{fact}

\begin{fact}[The chart form of the spin isomorphism]\label{c:mass.chart-form-spin-isomorphism}
In the conformal spinor chart, the dilation, the translations and the special conformal generators correspond to the grading matrix and to the upper and lower off-diagonal blocks $b_x$ and $c_x$.
\emph{Source:} T11, Lem 5.3 (\href{https://doi.org/10.5281/zenodo.22722565}{doi:10.5281/zenodo.22722565}).
\end{fact}

\begin{fact}[The Euclidean section, and its period]\label{c:singularity.euclidean-section-period}
The Euclidean section of the Schwarzschild field is smooth at its tip, the Euclidean counterpart of the bifurcation sphere, if and only if imaginary Killing time has period $p = 2\pi/\kappa = 8\pi M$; any other period leaves a conical point of total angle $\kappa p$.
\emph{Source:} T5, Prop 7.18 (\href{https://doi.org/10.5281/zenodo.22721790}{doi:10.5281/zenodo.22721790}).
\end{fact}

\begin{fact}[The chain]\label{c:standard-model.chain}
Under (S), the Pati-Salam chain with the framework's content unifies the three couplings at one loop with exact rational scales, and at two loops the version with a complex bidoublet and the $(15,2,2)$ passes the proton condition by a factor of six in $M_{U}$.
\emph{Source:} T12, Prop 10.10 (\href{https://doi.org/10.5281/zenodo.22722669}{doi:10.5281/zenodo.22722669}).
\end{fact}

\begin{fact}[The connection is forced]\label{c:standard-model.connection-forced}
A connection on $\mathcal{S} \otimes \mathcal{H}$ induced from a $\mathrm{Spin}(2,4)$-connection and a connection preserving the hierarchy's factorisation and Hermitian structure is exactly an $\so(2,4)$ part plus an $\su(2) \oplus \su(3) \oplus \uu(1)$ gauge field.
\emph{Source:} T12, Prop 5.7 (\href{https://doi.org/10.5281/zenodo.22722669}{doi:10.5281/zenodo.22722669}).
\end{fact}

\begin{fact}[The twistor form on the conformal spinor space is the grading of the central rotation $Z = e_{0}\nu$ against the Hilbert form of the ladder-pair module, so the twistor attachment's conformal energy $|s| + 1 + n$ is the nonnegative generator of $Z$; every future translation, null special conformal transformation and de Sitter horizon generator is nonnegative on its modules, and the attachment's Fock vacuum satisfies the positivity and invariance clause of (N$_{+}$) on the local horizons of Minkowski space (and of de Sitter space for conformally invariant fields); on those horizons (N$_{+}$) is the choice of that vacuum as the state of the light rays, a restatement of the premise and not a derivation of it]\label{c:synthesis.attachment-conformal-energy-is-generated-by-z}
On the conformal spinor space the invariant Hermitian form is, up to a real factor, $\ip{\psi}{iZ\phi}$ with $Z = e_{0}\nu$, so $Z$ grades the twistor form into the Hilbert form; quantised, $Z$ generates the twistor attachment's conformal energy $\tfrac12(N_{a} + N_{b}) + 1 = |s| + 1 + n$ with the stored multiplicities, future translations, null special conformal transformations and the de Sitter horizon generator are nonnegative with it, and the attachment's Fock vacuum satisfies the positivity and invariance clause of (N$_{+}$) on Minkowski local horizons (and de Sitter for conformal fields) with no thermal input; on those horizons (N$_{+}$) becomes the choice of that vacuum, a restatement and not a derivation of the premise; the eternal hole is not reached.
\emph{Source:} Transport Lab claim synthesis.attachment-conformal-energy-is-generated-by-z, version 1.
\end{fact}

\begin{fact}[On the attachment's helicity modules, the $L^{2}$ coefficient of the zero-point energy regulated by a smooth cut-off $L$ on conformal energy is $-(N_{V} - N_{W}/4)$ times one moment of the cut-off, reproducing the standard Sakharov weights; with $N_{V} = 12$, sixteen Weyl fields (right-handed neutrino included) in each of $N_{g}$ separate generations and conformally coupled scalars it vanishes exactly at $N_{g} = 3$, and a sharp cut-off at a floor adds terms that are not local]\label{c:synthesis.attachment-modules-give-sakharov-weights-cancelling-at-three-separate-generations}
On the attachment's helicity modules, the $L^{2}$ coefficient of the zero-point energy under a smooth cut-off on conformal energy (frame-dependent, defined on Weyl-flat fields) is $-(N_{V} - N_{W}/4)$ times one moment of the cut-off, the standard Sakharov weights; with $N_{V} = 12$, sixteen Weyl fields (right-handed neutrino included) in each of $N_{g}$ separate generations and conformal scalars it vanishes exactly at $N_{g} = 3$, which read on the Einstein static field is the induced Einstein-Hilbert coefficient if the regulated action is local and covariant (argued); a sharp floor cut-off adds non-local terms.
\emph{Source:} Transport Lab claim synthesis.attachment-modules-give-sakharov-weights-cancelling-at-three-separate-generations, version 2.
\end{fact}

\begin{fact}[For a hole formed by collapse, with the outgoing sector reduced to one chiral CAR field per light ray (J3), the ground-state condition of (N$_{+}$) posed in its short-distance form on the ingoing sheet through each point of the future horizon, whose generator has both halves, gives the late-time state of the near-horizon exterior outgoing modes as the continuum I-state at $\beta = 2\pi/\kappa$: in the Painlev\'e-Gullstrand gauge the Kruskal function $U = -e^{-\kappa t}(s-1)e^{s^{2}/2+s}$ is a regular parameter on that generator, dilated at rate $\kappa$ by the symmetry, so the late-time limit is exactly the scaling limit, which is the ground state for every kernel of the form ground-state singularity plus a remainder bounded near the horizon point, a class the same in every regular parameter for the weight-$\tfrac12$ kernel (sufficient condition; the null-shell realisation and the forcing of the temperature are a companion result)]\label{c:synthesis.collapsing-hole-flux-from-ground-state-on-the-ingoing-sheet}
For a collapse hole, with the outgoing sector reduced to one chiral CAR field per light ray (J3), the thermal condition is (N$_{+}$) posed on the ingoing sheet through each point of the future horizon, whose generator has both halves, in its local (scaling-limit) form; in the Painlev\'e-Gullstrand gauge $U = -e^{-\kappa t}(s-1)e^{s^{2}/2+s}$ is a regular parameter on that generator dilated at rate $\kappa$ by the symmetry, the late-time state of the near-horizon exterior outgoing modes is exactly the scaling limit and hence the continuum I-state at $\beta = 2\pi/\kappa$ (Hawking's flux at infinity under Premise (I)); the condition is sufficient, holding for every state whose generator kernel is the ground-state singularity plus a bounded remainder, a class the same in every regular parameter for the weight-$\tfrac12$ kernel (Schwarzian difference), and the Boulware-type state is outside it.
\emph{Source:} Transport Lab claim synthesis.collapsing-hole-flux-from-ground-state-on-the-ingoing-sheet, version 1.
\end{fact}

\begin{fact}[The cosmological horizon's boost is the frozen point's time translation]\label{c:synthesis.de-sitter-static-boost-frozen-point}
The static time of de Sitter space is the bivector $B_{\Lambda} = e_{0}\wedge W$ in the frozen point's stabiliser, with $B_{\Lambda}^{2} = \kappa_{\Lambda}^{2} = Q(I_{\Lambda}) = \Lambda/3$: the cosmological horizon's surface gravity squared is the frozen point's norm.
\emph{Source:} Transport Lab claim synthesis.de-sitter-static-boost-frozen-point, version 5.
\end{fact}

\begin{fact}[The standard null-cone form of the Einstein condition, translated into vacuum language: infinity vacua unexcited along their own light rays]\label{c:synthesis.einstein-condition-is-infinity-vacuum-incidence}
A translation, not a derivation: the standard equivalence ``Einstein if and only if $\Sch(\ell,\ell) = 0$ for all null $\ell$'' reads, through the Fock vacua annihilated by the pure-scale tractor, as ``every infinity vacuum is unexcited along its own light ray'' ($\Ric = (R/4)g$, $R$ then constant by the contracted Bianchi identity), with ``unexcited in every direction'' giving $\Ric = 0$, one unexcited vacuum giving $\Ric = 2c\,\ell\otimes\ell$, and, for a conserved source $T$, own-ray classes $4\pi G\,T(\ell,\ell)$ giving the Einstein equations with $\Lambda$ an integration constant; point vacua are unexcited along their own rays in every metric.
\emph{Source:} Transport Lab claim synthesis.einstein-condition-is-infinity-vacuum-incidence, version 3.
\end{fact}

\begin{fact}[The Euler part of the series' horizon entropy on a patch of a cut is an edge term, $\alpha$ times the holonomy angle of the cut's own rotation-gauge field around the patch edge; with the edge's geodesic-curvature term (the two-dimensional Gibbons-Hawking term, a choice within the edge ambiguity of the Noether charge, argued not derived) it becomes the constant $2\pi\alpha\chi(P) = -3\pi\chi(P)/2\Lambda$, so with that edge term the Wald reading of the series' action and the Clausius route agree on every entropy change of an open patch and differ by a topological constant that no additive edge term built from the field and the cut removes, and which for $\Lambda > 0$ makes the entropy negative on every disk patch of area below $6\pi/\Lambda$]\label{c:synthesis.euler-entropy-of-a-horizon-patch-is-an-edge-term}
The Euler part of the series' Jacobson-Myers entropy of a horizon patch is exact on the cut, $\alpha\oint_{\partial P}\omega_{\Sigma}$ (the holonomy angle of the cut's own rotation-gauge field around the edge; read as the wall counterterm one dimension down); with the edge's geodesic-curvature term (the two-dimensional Gibbons-Hawking term, a choice within the edge ambiguity of the Noether charge, argued not derived) it becomes $2\pi\alpha\chi(P) = -3\pi\chi(P)/2\Lambda$, so with that edge term the action and Clausius routes agree on every entropy change of an open patch (zero change on radial Schwarzschild-de Sitter generators) and differ only by that topological constant, which no additive edge term built from the field and the cut removes and which for $\Lambda > 0$ makes the entropy negative on every disk patch of area below $6\pi/\Lambda$ and which the stored normalisation (Z) or the integrated reading of (Loc) excludes. Transport-absent (the series supplies $\alpha$).
\emph{Source:} Transport Lab claim synthesis.euler-entropy-of-a-horizon-patch-is-an-edge-term, version 2.
\end{fact}

\begin{fact}[The tractor connection as a rotation-gauge field of the conformal model on flat spacetime]\label{c:synthesis.gtg-tractor-gauge-field-reproduces-tractor-connection}
On flat spacetime the tractor connection is the $\Cl(2,4)$ gauge field $\Omega_{T}(a) = \Omega(a) - \nInf \wedge \hb^{-1}(a) + \nO \wedge \Sch(\hb^{-1}(a))$: it reproduces the stored connection, its parallel scale tractor is the Einstein condition with $Q(I) = \Lambda/3$, and its algebraic field strength is minus the Weyl curvature.
\emph{Source:} Transport Lab claim synthesis.gtg-tractor-gauge-field-reproduces-tractor-connection, version 1.
\end{fact}

\begin{fact}[Imported: on every gravitational field with a globally hyperbolic derived metric the free Dirac field exists as field operators with a causal anticommutator, and its CAR algebras of local observables satisfy a generalised form of the Haag-Kastler axioms; the causal propagator reproduces every solution from its Cauchy data, and the anticommutator pairing is the conserved Cauchy inner product of the solutions the test spinors define (Dimock 1982)]\label{c:synthesis.imported-dimock-dirac-car-algebra-on-a-gravitational-field}
On every gravitational field with a globally hyperbolic derived metric the free Dirac field exists, with a causal anticommutator and CAR algebras of local observables satisfying a generalised form of the Haag-Kastler axioms; the causal propagator reproduces every solution from its Cauchy data and the anticommutator pairing is the conserved Cauchy inner product of the solutions; no state is selected (Dimock 1982).
\emph{Source:} J. Dimock, Dirac quantum fields on a manifold, Trans. Amer. Math. Soc. 269 (1982), 133--147., Sec. II.A (S = S+ - S-, after Thm. 2.1); Sec. II.B (Prop. 2.2, Thm. 2.3, Prop. 2.4(b)-(d)); Sec. III.A (the inner product on Cauchy data); Sec. III.B and Thm. 3.1 (the anticommutator).
\end{fact}

\begin{fact}[Imported: for the Dirac field the Hadamard form is equivalent to the wave front set condition and propagates from a neighbourhood of one Cauchy surface to the whole field (Kratzert 2000; Hollands 2001; Sahlmann and Verch 2001), on Robertson-Walker fields the Dirac states adiabatic of infinite order are Hadamard (Hollands 2001), and quasi-free Hadamard states have as normal-coordinate scaling limit the flat massless Dirac two-point function (Sahlmann and Verch 2001, Prop. 5.6, posed for Majorana spinors)]\label{c:synthesis.imported-dirac-hadamard-form-propagates-and-adiabatic-states-are-hadamard}
For the free Dirac field on a globally hyperbolic gravitational field, the Hadamard form is equivalent to the wave front set condition on both orderings and propagates from a neighbourhood of one Cauchy surface to the whole field (Kratzert 2000; Hollands 2001; Sahlmann and Verch 2001); on Robertson-Walker fields homogeneous Dirac states adiabatic of infinite order are Hadamard (Hollands 2001; the converse is not imported); quasi-free Hadamard states scale in normal coordinates to the flat massless Dirac two-point function (Sahlmann and Verch 2001, Prop. 5.6, posed for Majorana spinors; Sanders 2010 for the Dirac equivalence).
\emph{Source:} K. Kratzert, Singularity structure of the two point function of the free Dirac field on a globally hyperbolic spacetime, Ann. Phys. (Leipzig) 9 (2000), 475--498; arXiv:math-ph/0003015.; S. Hollands, The Hadamard condition for Dirac fields and adiabatic states on Robertson-Walker spacetimes, Commun. Math. Phys. 216 (2001), 635--661; arXiv:gr-qc/9906076.; H. Sahlmann and R. Verch, Microlocal spectrum condition and Hadamard form for vector-valued quantum fields in curved spacetime, Rev. Math. Phys. 13 (2001), 1203--1246; arXiv:math-ph/0008029., Sec. 5.2; Thm. 5.8; Def. 5.3; Prop. 5.6; Lemma 5.7; K. Sanders, The locally covariant Dirac field, Rev. Math. Phys. 22 (2010), 381--430; arXiv:0911.1304..
\end{fact}

\begin{fact}[Imported: the Einstein-Cartan matter coupling, in Gauge Theory Gravity form (Kibble 1961, Sciama 1964, Hehl, von der Heyde, Kerlick and Nester 1976, Lasenby, Doran and Gull 1998)]\label{c:synthesis.imported-einstein-cartan-matter-coupling}
With frame and connection varied independently and matter coupled minimally, the frame equation is $\mathcal{G}(a) + \Lambda a + 8\pi G\,T(a) = 0$ with $T$ the matter's energy-momentum, and the connection equation ties the torsion algebraically to the matter's spin density, vanishing without spin (Kibble 1961, Sciama 1964, Hehl et al. 1976; in Gauge Theory Gravity Lasenby, Doran and Gull 1998).
\emph{Source:} T. W. B. Kibble, Lorentz invariance and the gravitational field, J. Math. Phys. 2 (1961), 212--221.; D. W. Sciama, The physical structure of general relativity, Rev. Mod. Phys. 36 (1964), 463--469.; F. W. Hehl, P. von der Heyde, G. D. Kerlick and J. M. Nester, General relativity with spin and torsion: foundations and prospects, Rev. Mod. Phys. 48 (1976), 393--416.; A. N. Lasenby, C. J. L. Doran and S. F. Gull, Gravity, gauge theories and geometric algebra, Phil. Trans. R. Soc. Lond. A 356 (1998), 487--582..
\end{fact}

\begin{fact}[Imported: the massless two-point function $D_{+}$, real and even at spacelike separation; Pauli's free-field spin-statistics argument (1940); and the algebraic spin-statistics theorem of Guido and Longo (1995): for a modular covariant net of field algebras with the Reeh-Schlieder property and normal commutation relations the statistics grading is $\Gamma = U(2\pi)$]\label{c:synthesis.imported-free-field-spin-statistics-and-guido-longo-theorem}
$D_{+}(\xi) = -1/(4\pi^{2}(Q(\xi) - i0\,\xi_{0}))$, real and even at spacelike $\xi$; Pauli's free-field spin-statistics argument; and Guido-Longo: modular covariance, Reeh-Schlieder and normal commutation relations give $\Gamma = U(2\pi)$ with positive energy derived.
\emph{Source:} W. Pauli, The connection between spin and statistics, Phys. Rev. 58 (1940), 716--722.; D. Guido and R. Longo, An algebraic spin and statistics theorem, Commun. Math. Phys. 172 (1995), 517, arXiv:funct-an/9406005, Theorem 2.11 and Proposition 2.4.; A. Shalabi, M. H. Zambianco and T. R. Perche, The effect of curvature on local observables in quantum field theory, arXiv:2412.12294, equation (11).; R. F. Streater and A. S. Wightman, PCT, Spin and Statistics, and All That, Benjamin (1964)..
\end{fact}

\begin{fact}[Imported: the dictionary between Gauge Theory Gravity and curved-spacetime language]\label{c:synthesis.imported-gtg-dictionary}
On the invertible region, tangent vectors are $\hb^{-1}(a)$, the Levi-Civita derivative is $D_a$, geodesics have $D$-parallel velocity, and the Riemann, Ricci, Schouten and Weyl tensors and Killing fields are read from the field strength $\mathcal{R}(B)$ and $D$.
\emph{Source:} A. Lasenby, C. Doran and S. Gull, Gravity, gauge theories and geometric algebra, Phil. Trans. R. Soc. Lond. A 356 (1998), 487--582., \S\S 3.5, 4.3--4.4, 5.1--5.2, App.~C; C. Doran and A. Lasenby, Geometric Algebra for Physicists, Cambridge University Press (2003)., \S\S~13.5--13.6, 14.1; A. N. Lasenby, Geometric algebra, gravity and gravitational waves, Adv. Appl. Clifford Algebras (2019)., \S\S~4--5, eq.~(8.5) (arXiv:1912.05960); R. M. Wald, General Relativity, University of Chicago Press (1984)., \S~12.5.
\end{fact}

\begin{fact}[Imported: the local normality principle of algebraic quantum field theory (Haag), and local quasi-equivalence in a gravitational background (Haag, Narnhofer and Stein 1984)]\label{c:synthesis.imported-haag-local-normality}
Physical states are locally normal to the vacuum, quasi-equivalent to it on the algebra of every bounded region (Haag); in a gravitational background any two physical states are locally quasi-equivalent, and with a short-distance scaling-limit condition at a horizon this gives the Hawking temperature (Haag, Narnhofer and Stein 1984).
\emph{Source:} R. Haag, Local Quantum Physics: Fields, Particles, Algebras, 2nd edition, Springer, Berlin, 1996.; R. Haag, H. Narnhofer and U. Stein, On quantum field theory in gravitational background, Commun. Math. Phys. 94 (1984), 219--238..
\end{fact}

\begin{fact}[Imported: the Hadamard parametrix of a normally hyperbolic operator on a vector bundle has leading coefficient the square root of the Van Vleck determinant times the parallel transport of the operator's connection (Friedlander 1975; Bar, Ginoux and Pfaffle 2007, Ch. 2); for the squared Dirac operator that connection is the spin connection (Lichnerowicz), so U = Delta\^{}(1/2) R; the Hadamard form does not depend on the time function of its i0 prescription (Kay and Wald 1991; Sahlmann and Verch 2001, Sec. 5.1); and in four dimensions the retarded and advanced fundamental solutions on a causal domain are (1/4 pi) theta(+-(t - t')) [U delta(sigma) + V theta(sigma)] with V smooth (Friedlander 1975; Bar, Ginoux and Pfaffle 2007, Ch. 2)]\label{c:synthesis.imported-hadamard-parametrix-of-the-squared-dirac-operator-leads-with-van-vleck-root-times-spin-transport}
For a normally hyperbolic operator on a vector bundle the leading Hadamard coefficient is $\Delta^{1/2}$ times the parallel transport of the operator's connection (Friedlander 1975; Bar, Ginoux and Pfaffle 2007); for the squared Dirac operator that is the spin connection, so $U = \Delta^{1/2}R$ with $\dot R = -\tfrac12\Omega(\dot\gamma)R$; the Hadamard form does not depend on the time function of its $i0$ prescription (Kay and Wald 1991; Sahlmann and Verch 2001, Sec. 5.1); in four dimensions the retarded and advanced fundamental solutions on a causal domain are $\tfrac{1}{4\pi}\theta(\pm(t - t'))[U\delta(\sigma) + V\theta(\sigma)]$, $V$ smooth.
\emph{Source:} F. G. Friedlander, The Wave Equation on a Curved Space-Time, Cambridge University Press, Cambridge, 1975., the fundamental solutions of normally hyperbolic operators and their Hadamard form; cited from memory, not read; C. Bar, N. Ginoux and F. Pfaffle, Wave Equations on Lorentzian Manifolds and Quantization, ESI Lectures in Mathematics and Physics, European Mathematical Society, 2007; arXiv:0806.1036., Ch. 2 (Def. 2.2.1, Prop. 2.3.1; the formal and local fundamental solutions as sums of Hadamard coefficients times Riesz distributions); Sec. 1.3 (the function mu\_x); E. Poisson, The motion of point particles in curved spacetime, Living Rev. Relativity 7 (2004), 6; arXiv:gr-qc/0306052., Sec. 4.3.2, Eqs. (266), (268), (269), (273) (scalar Green functions); Sec. 4.4.2 (electromagnetic Green functions, U = parallel propagator times Delta\^{}(1/2)); H. Sahlmann and R. Verch, Microlocal spectrum condition and Hadamard form for vector-valued quantum fields in curved spacetime, Rev. Math. Phys. 13 (2001), 1203--1246; arXiv:math-ph/0008029., Sec. 3.5; Sec. 5.1; B. S. Kay and R. M. Wald, Theorems on the uniqueness and thermal properties of stationary, nonsingular, quasifree states on spacetimes with a bifurcate Killing horizon, Phys. Rep. 207 (1991), 49--136.; M. J. Radzikowski, Micro-local approach to the Hadamard condition in quantum field theory on curved space-time, Commun. Math. Phys. 179 (1996), 529--553.; L. Ronge, Formulas for Hadamard coefficients in terms of Green's operators, arXiv:2604.25763 (2026)., Def./Prop. 2.9, 2.10; Sec. 2.1; A. G. M. Lewis, Hadamard renormalization of a 2-dimensional Dirac field, Phys. Rev. D 101 (2020), 125019; arXiv:1910.05663., Eqs. (14), (26), (52).
\end{fact}

\begin{fact}[Imported: on an arbitrary four-dimensional spacetime, near the vertex of the future light cone of a point (inside a geodesically convex, globally hyperbolic neighbourhood), the trace of a Dirac field on the cone, its two transverse spinor components, is an L2 isometry from data on a spacelike slice (the flux identity of the conserved current), the characteristic Cauchy problem on the cone is uniquely solvable for L2 data, and the trace is an isomorphism of H1 onto the space of cone data (Haefner and Nicolas 2011, Theorems 1 and 2)]\label{c:synthesis.imported-haefner-nicolas-dirac-characteristic-problem-on-light-cones-near-the-vertex}
For the Dirac field of any mass on an arbitrary four-dimensional spacetime, near the vertex of the future light cone of a point (in a geodesically convex, globally hyperbolic neighbourhood, below a level $t = T$), the trace on the cone, two transverse spinor components, is an $L^{2}$ isometry from data on a spacelike slice by the flux identity of the conserved current, the characteristic problem is uniquely solvable for $L^{2}$ data, so the trace is unitary, and it is an isomorphism of $H^{1}$ onto the cone data (Häfner and Nicolas 2011, Theorems 1 and 2).
\emph{Source:} D. Häfner and J.-P. Nicolas, The characteristic Cauchy problem for Dirac fields on curved backgrounds, J. Hyperbolic Differ. Equ. 8 (2011), no. 3, 437--483; arXiv:0903.0515., Abstract; Sec. 3 (Def. 3.6, the domain D\_T); Sec. 4.1 (Def. 4.2, Eq. (14), Thm. 1); Sec. 4.4 (Def. 4.11, Thm. 2).
\end{fact}

\begin{fact}[Imported: Einstein's equation as an equation of state (Jacobson 1995)]\label{c:synthesis.imported-jacobson-clausius-1995}
If every local Rindler horizon has the Unruh temperature, entropy $\eta$ per area and obeys $\delta Q = T\,\delta S$ for the matter's boost-energy flux, then $G_{ab} + \Lambda g_{ab} = (2\pi/\eta)T_{ab}$ with $\Lambda$ undetermined (Jacobson 1995), the field equation $\mathcal{G}(a) + \Lambda a = -(2\pi/\eta)T(a)$ of the gravitational field in the store's sign.
\emph{Source:} T. Jacobson, Thermodynamics of spacetime: the Einstein equation of state, Phys. Rev. Lett. 75 (1995), 1260--1263..
\end{fact}

\begin{fact}[Imported: on a spacetime with a bifurcate Killing horizon there is at most one stationary quasi-free Hadamard state, KMS at the Hawking temperature on a wedge; none exists on extended Schwarzschild-de Sitter, and on Kerr none is argued (Kay and Wald 1991, with Kay and Lupo 2016)]\label{c:synthesis.imported-kay-wald-uniqueness-and-no-go}
On a spacetime with a bifurcate Killing horizon there is at most one stationary quasi-free Hadamard state, KMS at $\kappa/2\pi$ on a wedge (Kay and Wald 1991); none exists on extended Schwarzschild-de Sitter, and on Kerr none is argued for the scalar field; a gap in the proofs was filled by Kay and Lupo (2016, App. B).
\emph{Source:} B. S. Kay and R. M. Wald, Theorems on the uniqueness and thermal properties of stationary, nonsingular, quasifree states on spacetimes with a bifurcate Killing horizon, Phys. Rep. 207 (1991), 49--136.; B. S. Kay and U. Lupo, Non-existence of isometry-invariant Hadamard states for a Kruskal black hole in a box and for massless fields on 1+1 Minkowski spacetime with a uniformly accelerating mirror, Class. Quantum Grav. 33 (2016), 215001; arXiv:1502.06582., App. B; footnotes 2 and 4; B. S. Kay, Quantum field theory in curved spacetime (2nd edition), in Encyclopedia of Mathematical Physics, 2nd edition (2024); arXiv:2308.14517..
\end{fact}

\begin{fact}[Imported: quasi-free Hadamard states are locally quasi-equivalent, for the Klein-Gordon field (Verch 1994) and for the Dirac field (D'Antoni and Hollands 2006)]\label{c:synthesis.imported-local-quasi-equivalence-of-hadamard-states}
Any two quasi-free Hadamard states are locally quasi-equivalent on the algebra of every open relatively compact region, for the Klein-Gordon field (Verch 1994) and for the Dirac field, there with local unitary equivalence and type III$_{1}$ local algebras (D'Antoni and Hollands 2006).
\emph{Source:} R. Verch, Local definiteness, primarity and quasiequivalence of quasifree Hadamard quantum states in curved spacetime, Commun. Math. Phys. 160 (1994), 507--536.; C. D'Antoni and S. Hollands, Nuclearity, local quasiequivalence and split property for Dirac quantum fields in curved spacetime, Commun. Math. Phys. 261 (2006), 133--159; arXiv:math-ph/0106028., Thms. III.1, III.2.
\end{fact}

\begin{fact}[Imported: in four dimensions the natural second-order field equations of a metric are the Einstein map plus a constant (Lovelock), and the natural Lagrangians with no field equations are the Euler and Pontryagin densities up to divergences (Anderson and Duchamp; Anderson)]\label{c:synthesis.imported-lovelock-and-anderson-duchamp-natural-lagrangians}
In four dimensions a natural divergence-free second-order symmetric expression in the metric is $c_{1}\mathcal{G} + c_{2}g$ (Lovelock 1971, 1972), and a natural Lagrangian with identically vanishing Euler-Lagrange expression is $c_{E}E_{4} + c_{P}P_{4}$ plus a divergence (Anderson and Duchamp 1980; Anderson 1984).
\emph{Source:} D. Lovelock, The Einstein tensor and its generalizations, J. Math. Phys. 12 (1971), 498--501.; D. Lovelock, The four-dimensionality of space and the Einstein tensor, J. Math. Phys. 13 (1972), 874--876.; I. M. Anderson and T. Duchamp, On the existence of global variational principles, Amer. J. Math. 102 (1980), 781--868.; I. M. Anderson, Natural variational principles on Riemannian structures, Ann. of Math. (2) 120 (1984), 329--370..
\end{fact}

\begin{fact}[Imported: the stress of an electromagnetic field $F$ is $T(a) = -\tfrac12 FaF$ (Hestenes; Doran and Lasenby)]\label{c:synthesis.imported-maxwell-stress-spacetime-algebra}
The electromagnetic stress in the spacetime algebra is $T(a) = -\tfrac12 FaF$: symmetric, trace-free, with energy density $\tfrac12(E^{2} + B^{2})$ and, for an axis-aligned field, components $\mathrm{diag}(\rho, -\rho, \rho, \rho)$ (Hestenes; Doran and Lasenby).
\emph{Source:} D. Hestenes, Space-Time Algebra, Gordon and Breach (1966).; C. Doran and A. Lasenby, Geometric Algebra for Physicists, Cambridge University Press (2003), Section 7.4..
\end{fact}

\begin{fact}[Imported: the microlocal form of the Hadamard condition (Radzikowski 1996), and for the Dirac field on both orderings of the two-point function (Hollands 2001), under which Hadamard states differ by smooth kernels]\label{c:synthesis.imported-microlocal-hadamard-condition}
A quasi-free Klein-Gordon state is Hadamard exactly when its two-point function satisfies the wave front set spectral condition $\mathrm{WF}'(\Lambda) = C^{+}$ (Radzikowski 1996); a quasi-free Dirac state is Hadamard when $\mathrm{WF}'(\mathcal{G}^{(\pm)}) = C^{(\pm)}$ on both orderings, imposed separately (Hollands 2001); Hadamard states differ by smooth kernels.
\emph{Source:} M. J. Radzikowski, Micro-local approach to the Hadamard condition in quantum field theory on curved space-time, Commun. Math. Phys. 179 (1996), 529--553.; S. Hollands, The Hadamard condition for Dirac fields and adiabatic states on Robertson-Walker spacetimes, Commun. Math. Phys. 216 (2001), 635--661; arXiv:gr-qc/9906076., Def. IV.1; Thm. IV.2; C. J. Fewster, Lectures on quantum energy inequalities, arXiv:1208.5399., Secs. 2.3, 3.2.
\end{fact}

\begin{fact}[Imported: Nordstr\"om's scalar theory of gravitation in the Einstein-Fokker form (vanishing Weyl curvature, scalar curvature proportional to the stress trace), which deflects no light and gives minus one sixth of Einstein's perihelion advance, so is excluded by observation]\label{c:synthesis.imported-nordstrom-scalar-theory-einstein-fokker-form}
Nordstr\"om's theory in the Einstein-Fokker form, $\mathcal{W} = 0$ and $\mathcal{R} = -24\pi G\tau$ in the store's sign ($R = 24\pi G T$ in the sources'), deflects no light ($\gamma = -1$) and gives $-\tfrac16$ of Einstein's perihelion advance, so it is excluded by observation (Einstein and Fokker 1914; Deruelle and Sasaki 2011).
\emph{Source:} A. Einstein and A. D. Fokker, Die Nordstr\"omsche Gravitationstheorie vom Standpunkt des absoluten Differentialkalk\"uls, Ann. Phys. (Leipzig) 349 (1914), 321--328.; G. Nordstr\"om, Zur Theorie der Gravitation vom Standpunkt des Relativit\"atsprinzips, Ann. Phys. (Leipzig) 42 (1913), 533--554.; N. Deruelle and M. Sasaki, Conformal transformations and Nordstr\"om's scalar theory of gravity, Prog. Theor. Phys. Suppl. 190 (2011), 143--154; arXiv:1012.5386., Sec. II.3; Sec. III.1, Eq. (3.7); C. W. Misner, K. S. Thorne and J. A. Wheeler, Gravitation, Freeman, San Francisco (1973)., Exercise 7.1.
\end{fact}

\begin{fact}[Imported: light-front quantisation of free fields: on a null plane the free fields of different masses are unitarily equivalent and generate an irreducible algebra (Leutwyler, Klauder and Streit 1970), and the restriction of a free field to a null surface carries no mass except in zero modes, the mass entering the Dirac field only through its constrained half (Wall 2012)]\label{c:synthesis.imported-null-plane-restriction-of-free-fields-is-mass-independent}
On a null plane free fields of different masses are unitarily equivalent and generate an irreducible algebra (Leutwyler, Klauder and Streit 1970); restricted to a null surface a free field's algebra and null energy carry no mass except in zero modes, the Dirac field's mass entering only its constrained half (Wall 2012).
\emph{Source:} H. Leutwyler, J. R. Klauder and L. Streit, Quantum field theory on lightlike slabs, Nuovo Cimento A 66 (1970), 536--554.; A. C. Wall, A proof of the generalized second law for rapidly changing fields and arbitrary horizon slices, Phys. Rev. D 85 (2012), 104049; arXiv:1105.3445., Secs. 3.2, 4.1.
\end{fact}

\begin{fact}[Imported: quasi-equivalence of quasi-free CAR states and the Powers-Stormer inequality (Powers-Stormer), and disjointness of KMS states at different temperatures (Takesaki)]\label{c:synthesis.imported-powers-stormer-takesaki-disjointness}
Gauge-invariant quasi-free CAR states $\omega_{A}$, $\omega_{B}$ are factor states, quasi-equivalent iff $A^{1/2} - B^{1/2}$ and $(1 - A)^{1/2} - (1 - B)^{1/2}$ are Hilbert-Schmidt and otherwise disjoint, and $\|A^{1/2} - B^{1/2}\|_{2}^{2} \le \|A - B\|_{1}$ for bounded positive $A$, $B$ (Powers-Størmer); KMS states at different temperatures are disjoint when one induces a type III algebra (Takesaki).
\emph{Source:} R. T. Powers and E. Stormer, Free states of the canonical anticommutation relations, Commun. Math. Phys. 16 (1970), 1--33.; M. Takesaki, Disjointness of the KMS-states of different temperatures, Commun. Math. Phys. 17 (1970), 33--41..
\end{fact}

\begin{fact}[Imported: the Hadamard form of a two-point function in a neighbourhood of one Cauchy surface propagates to the whole spacetime (Fulling, Sweeny and Wald 1978)]\label{c:synthesis.imported-propagation-of-hadamard-form}
For the free scalar field, a two-point function of Hadamard form in a neighbourhood of one Cauchy surface is of Hadamard form everywhere: Cauchy evolution preserves the Hadamard form (Fulling, Sweeny and Wald 1978).
\emph{Source:} S. A. Fulling, M. Sweeny and R. M. Wald, Singularity structure of the two-point function in quantum field theory in curved spacetime, Commun. Math. Phys. 63 (1978), 257--264..
\end{fact}

\begin{fact}[Imported: the focusing and shear equations for a null geodesic congruence (Raychaudhuri 1955, Sachs 1961)]\label{c:synthesis.imported-raychaudhuri-null-congruence}
For a null geodesic congruence of physical vectors, $(k \cdot D)\theta = -\theta^{2}/2 - \sigma^{2} + \omega^{2} - \Ric^{\circ}(k,k)$, with $-\Ric^{\circ}(k,k) = k \cdot \mathcal{R}(k)$ in the store's sign, and $d\sigma/d\lambda = -\theta\sigma - C(\cdot,k,\cdot,k)$ on the screen space (Raychaudhuri 1955, Sachs 1961).
\emph{Source:} A. Raychaudhuri, Relativistic cosmology. I, Phys. Rev. 98 (1955), 1123--1126.; R. K. Sachs, Gravitational waves in general relativity. VI. The outgoing radiation condition, Proc. R. Soc. Lond. A 264 (1961), 309--338.; R. M. Wald, General Relativity, University of Chicago Press (1984)..
\end{fact}

\begin{fact}[Imported: the one-loop induced Einstein-Hilbert coefficient of Sakharov's induced gravity, proportional to $N_{0} + N_{1/2} - 4N_{1}$ (Broda and Szanecki), with the curvature-linear heat-kernel density $\mathrm{tr}(E + R/6)$ (Vassilevich), the count $n_{1/2} = 4n_{1}$, $48 = 4 \times 12$, for the Standard Model with right-handed neutrinos (Boyle and Turok), and no Einstein-tensor ambiguity for massless fields (Tichy and Flanagan)]\label{c:synthesis.imported-sakharov-induced-gravity-coefficients}
At one loop the induced Einstein-Hilbert coefficient is proportional to $N_{0} + N_{1/2} - 4N_{1}$ (Broda and Szanecki; $N_{0}(1 - 6\xi)$ for curvature-coupled scalars, from the heat-kernel density $\mathrm{tr}(E + R/6)$, Vassilevich); the Standard Model with right-handed neutrinos has $n_{1/2} = 4n_{1} = 48$ (Boyle and Turok); massless fields have no Einstein-tensor ambiguity in the locally covariant stress tensor (Tichy and Flanagan).
\emph{Source:} D. V. Vassilevich, Heat kernel expansion: user's manual, Phys. Rep. 388 (2003), 279--360; arXiv:hep-th/0306138., Eq. (4.27); B. Broda and M. Szanecki, Induced gravity and gauge interactions revisited, Phys. Lett. B (2009); arXiv:0809.4203., Eqs. (5), (8); L. Boyle and N. Turok, Cancelling the vacuum energy and Weyl anomaly in the standard model with dimension-zero scalar fields, arXiv:2110.06258., Eq. (17); W. Tichy and E. E. Flanagan, How unique is the expected stress-energy tensor of a massive scalar field?, Phys. Rev. D 58 (1998), 124007; arXiv:gr-qc/9807015., Sec. I.A, Eqs. (1.3)--(1.5).
\end{fact}

\begin{fact}[Imported: the integrability condition of the twistor equation, $\Psi_{ABCD}\omega^{D} = 0$, restricts the Weyl curvature at every point where a nonzero solution exists to Petrov type N (with $\omega$ a repeated principal spinor) or O (Penrose; Penrose and Rindler; Gasperin and Pais)]\label{c:synthesis.imported-twistor-equation-integrability-petrov-n-or-o}
The integrability condition of the twistor equation, $\Psi_{ABCD}\omega^{D} = 0$, forces the Weyl curvature at every point with a nonzero solution to be of Petrov type N (with $\omega$ the repeated principal spinor) or O (Penrose and Rindler 1986; Gasperin and Pais 2024).
\emph{Source:} R. Penrose and W. Rindler, Spinors and Space-Time, vol. 2: Spinor and Twistor Methods in Space-Time Geometry, Cambridge University Press (1986)., Ch. 6; E. Gasperin and F. Pais, Twistor initial data characterisation of pp-waves, arXiv:2404.05883..
\end{fact}

\begin{fact}[On the integrated reading of (Loc), the Clausius relation with the de Sitter calibration forces a continuous entropy density to 1/4 on every admissible field whose local data approach calibrated data; so a continuous density depending on curvature alone that keeps the flat Friedmann fields with $w \in [0,1]$ and $\Lambda \ge 0$ admissible is 1/4 at every curvature they pass through, the Einstein static universe's among them; a density that depends on matter as well escapes with every field satisfying Einstein's equations at the calibrated coupling admissible]\label{c:synthesis.integrated-loc-curvature-clause-forced-matter-clause-independent}
On the integrated reading, Clausius plus the de Sitter calibration forces a continuous horizon entropy density to 1/4 on every admissible field whose local data approach de Sitter or flat data; a continuous density depending on curvature alone that keeps the flat Friedmann fields with $w \in [0,1]$, $\Lambda \ge 0$ admissible is therefore 1/4 at every isotropic Weyl-free curvature with $\beta > 0$, $-2\beta \le \alpha < \beta$ they pass through, the Einstein static universe's included; a density depending on matter, $\tfrac14(1 + cQ)$ with $Q$ the square of the departure from Einstein's equations, keeps every Einstein solution admissible and is $\tfrac12$ on an Einstein static universe.
\emph{Source:} Transport Lab claim synthesis.integrated-loc-curvature-clause-forced-matter-clause-independent, version 1.
\end{fact}

\begin{fact}[Under the Gupta-Bleuler form of the Krein premise, the Krein quartet of a covariantly gauge-fixed vector drops level by level from the induced Einstein-Hilbert coefficient under an energy cut-off on the Einstein static field, but under a heat-kernel cut-off on the gauge-fixed operators the vector's weight is $\alpha - 5$, equal to the two-helicity value $-4$ only at the Gupta-Bleuler value $\alpha = 1$; so under that cut-off the series' content induces $12(\alpha - 1) + 8(1 - 6\xi_{R})$ minimal-scalar units, fixed by the content only together with the regulator and the gauge]\label{c:synthesis.krein-quartet-drops-from-the-induced-einstein-hilbert-coefficient-only-under-a-regulator-that-treats-it-alike}
Under the Gupta-Bleuler form of the Krein premise, the unphysical quartet of a gauge vector drops level by level from the induced Einstein-Hilbert coefficient under an energy cut-off on the Einstein static field, but under a heat-kernel cut-off on the gauge-fixed operators the vector's weight is $\alpha - 5$ ($-4$ only at the Gupta-Bleuler value $\alpha = 1$), so the series' content induces $12(\alpha - 1) + 8(1 - 6\xi_{R})$ minimal-scalar units under the covariant cut-off and $8(1 - 6\xi_{R})$ under the energy cut-off.
\emph{Source:} Transport Lab claim synthesis.krein-quartet-drops-from-the-induced-einstein-hilbert-coefficient-only-under-a-regulator-that-treats-it-alike, version 2.
\end{fact}

\begin{fact}[On a horizon's null plane the canonical part of a spacetime-algebra spinor is the image of Lambda(m), two complex channels, one per chirality, with helicity about the generators equal to minus the chirality; so, given one (Q) version 3 array per complex channel, one massless Dirac species supplies two arrays per light ray (ln 2 at the Hawking temperature) and one Weyl field one]\label{c:synthesis.light-ray-good-part-of-a-spinor-is-two-chirality-arrays}
In $\Cl(1,3)$ the flux of the Dirac current through a horizon's null plane $N_{k}$ is $2|\Lambda(m)\psi|^{2}$, so the canonical part there is the image of $\Lambda(m) = m\gamma_{0}/2$, real dimension $4$, of boost weight $\tfrac12$ in the carrier's convention, split by chirality into two complex channels with helicity about $k$ minus the chirality; given one (Q) version 3 array per channel, one massless Dirac species gives two arrays per light ray, $\ln 2$ at the Hawking temperature, a Weyl field one, and the per-light-ray total over fields counts two arrays per registered unit.
\emph{Source:} Transport Lab claim synthesis.light-ray-good-part-of-a-spinor-is-two-chirality-arrays, version 2.
\end{fact}

\begin{fact}[Quantised as an elementary field with $i = I_{4}$, the rung field is linear on all eight real components of the complexified rung summand, and on that summand (on values, not on spacetime points) the hypercharge rotation is the rung boost at $I_{4}$-imaginary rapidity; besides the stored neutral mode and charged pair it carries a neutral normal pair, the directions that change $Q$, of mass squared $2W' + 4N(k_{10} \pm r)$, and the null-lift constraint is the limit $k_{10} \to \infty$; under the Gildener-Weinberg formula and decoupling the normal pair enters the scalon sum (the standard two-doublet sum rule), so the charged pair lies below $456$ GeV (below $383$ GeV when the normal pair is the heavier), the normal pair below $542$ GeV, and $h \to \gamma\gamma$ below $0.900$ of the Standard Model rate for every charged mass between $m_{h}/2$ and $456$ GeV]\label{c:synthesis.linear-rung-field-has-a-neutral-normal-pair-in-the-scalon-sum}
With $i = I_{4}$ the elementary rung field is linear on eight components (on the rung summand, hypercharge is the rung boost at $I_{4}$-imaginary rapidity) and adds a neutral normal pair at $2W' + 4N(k_{10} \pm r)$ to the stored spectrum, the constraint being $k_{10} \to \infty$; under Gildener-Weinberg and decoupling $2m_{\pm}^{4} + m_{H}^{4} + m_{A}^{4} \in (2 \cdot 450^{4}, 2 \cdot 456^{4})$ GeV$^{4}$, so $m_{\pm} < 456$ GeV, $m_{H,A} < 542$ GeV, and $\mu_{\gamma\gamma} < 0.900$ for every charged mass in $(m_{h}/2, 456]$ GeV. (Excluded on 2 November 2026 under the symmetric coupling reading: Particle Data Group 2026 $\mu_{\gamma\gamma} = 1.10 \pm 0.06$, $B$ mixing over $50$ times the measured rate, quartics at $4\pi$ by $0.8$ to $4$ TeV; the level branch is the working branch.)
\emph{Source:} Transport Lab claim synthesis.linear-rung-field-has-a-neutral-normal-pair-in-the-scalon-sum, version 3.
\end{fact}

\begin{fact}[The curvature-independence clause (Loc) is independent of the Clausius relation, of (R$_{\mathrm{loc}}$) and of every $\Lambda$-vacuum horizon datum of the store; a density $\tfrac14(1 + c\,S_{ab}S^{ab})$ passes all of them and differs from $\tfrac14$ in the Einstein static universe, and on the integrated reading such failures are confined to spacetimes with no calibrated vacuum point]\label{c:synthesis.loc-independent-of-clausius-and-killing-data}
A local-horizon density $\tfrac14(1 + c\,S_{ab}S^{ab})$ equals $\tfrac14$ on every $\Lambda$-vacuum, so it keeps every $\Lambda$-vacuum horizon datum, (R$_{\mathrm{loc}}$) and the Clausius relation, yet is $\tfrac14(1 + 3c/a^{4})$ in the dust-filled Einstein static universe where Clausius still balances: (Loc) is independent of them; on the integrated reading such failures need a spacetime with no calibrated vacuum point.
\emph{Source:} Transport Lab claim synthesis.loc-independent-of-clausius-and-killing-data, version 4.
\end{fact}

\begin{fact}[On local Rindler horizons and small diamonds, clauses (ii) and (iii) of (C$_{\mathrm{loc}}$), with (iii) read on the patch's whole mode algebra and the boost acting on the modes by phases, are one condition, joint KMS for the local boost at rapidity period $2\pi$, the period at which the local boost's Euclidean section is regular at leading order]\label{c:synthesis.local-boost-temperature-is-local-regularity}
For any patch carrying modes on which the local boost acts by phases (assumed), clause (iii) of (C$_{\mathrm{loc}}$), read on the patch's whole mode algebra, is equivalent to clause (ii) (KMS at $\beta$ iff faithful with modular group the boost at $-\beta s$, on every window), and both together are (R$_{\mathrm{loc}}$), the joint KMS condition of the I-state for the local boost at rapidity period $2\pi$, a new premise stronger than the local form of (R); derived is only the geometric half, that this period is where the local boost's Euclidean section is regular for every field strength at leading order, giving $T = a/2\pi$ up to a relative correction $-\tfrac32k\rho^{2} + O(\rho^{3})$.
\emph{Source:} Transport Lab claim synthesis.local-boost-temperature-is-local-regularity, version 3.
\end{fact}

\begin{fact}[On a horizon generator the weight-$\tfrac12$ scaling preserves the Hilbert-Schmidt norm of a kernel shrunk onto a point, so a gauge-invariant quasi-free state locally normal to the null vacuum at a point is locally the null vacuum there in every regular parameter, and a stationary gauge-invariant quasi-free state locally normal to it at the edge is the null vacuum of (N$_{+}$), positivity included; every other stationary gauge-invariant quasi-free state, the half-thermal products at every period included, is locally disjoint from it at the edge]\label{c:synthesis.local-normality-at-the-edge-gives-the-null-ground-state}
On a horizon generator the weight-$\tfrac12$ scaling preserves the Hilbert-Schmidt norm of a kernel shrunk onto a point, so for gauge-invariant quasi-free states local normality to the null vacuum at a point (via Powers-Stormer, a square-integrable difference) gives the local form of (N$_{+}$) there in any regular parameter, and a stationary such state locally normal to it at the edge is the null vacuum of (N$_{+}$), positivity included; every other stationary gauge-invariant quasi-free state differs by a degree $-1$ kernel that is not square integrable at the edge, so the half-thermal products, $c = 1$ included, are locally disjoint from the vacuum there.
\emph{Source:} Transport Lab claim synthesis.local-normality-at-the-edge-gives-the-null-ground-state, version 2.
\end{fact}

\begin{fact}[The imported MacDowell-Mansouri functional, in the series' frozen-point form, is a dynamical principle for the whole gauge field that fixes the trace-free part: its frame equation is twice the dual of $\mathcal{G}(a) + \Lambda a$, ten conditions with the Weyl part dropping out, and its restriction to pure dilations is the single Yamabe trace; with the standard (imported) Einstein-Cartan matter coupling its on-shell Schouten tensor is the stored derivative of the scale tractor, so the infinity-vacuum incidence is its vacuum case]\label{c:synthesis.macdowell-mansouri-principle-fixes-the-trace-free-part}
The imported MacDowell-Mansouri functional, in the series' frozen-point form $S \sim Q(I)^{-1}\int\epsilon(I, V, F, F)$, varies the whole field $(\hb, \Omega)$; its frame equation is $E(a) = 2(\mathcal{G}(a) + \Lambda a)I$ (components: $\epsilon_{abcd}e^{b}\wedge F^{cd}$), rank $10$ on the curvature with kernel the Weyl part, so it fixes the trace-free part, while its restriction to pure dilations is the single Yamabe trace of the scale-only no-go; with the standard (imported) Einstein-Cartan matter coupling and an independent source it gives clause (7), on shell $\Sch = \tfrac{\Lambda}{6}g - 4\pi G(T - \tfrac{\tau}{3}g)$, the stored derivative of the scale tractor, and in vacuum the infinity-vacuum incidence. It settles the search for such a principle only given this imported action principle (Einstein's equations as a variational premise); transport-absent, and its Wald entropy carries the Euler term that (Loc) excludes.
\emph{Source:} Transport Lab claim synthesis.macdowell-mansouri-principle-fixes-the-trace-free-part, version 2.
\end{fact}

\begin{fact}[On a mode space where the boost acts by phases, it dilates no translation group, so the Borchers route to (R) needs the continuum]\label{c:synthesis.point-spectrum-boost-admits-no-half-sided-translations}
A boost that acts by phases on a generating set of modes (pure point spectrum: windows, the countable mode space, the array of (Q)) dilates no strongly continuous translation group, even on the GNS space of a KMS state, so the Borchers and Bisognano-Wichmann route to (R) must be posed on the continuum mode space.
\emph{Source:} Transport Lab claim synthesis.point-spectrum-boost-admits-no-half-sided-translations, version 3.
\end{fact}

\begin{fact}[(R$_{\mathrm{loc}}$), with the local-horizon modes labelled by null direction, is independent of the present premises, and what it adds off Killing horizons is the matter coupling; a local rapidity period $2\pi c$ gives Einstein's equations with $G_{\mathrm{cl}} = cG$]\label{c:synthesis.r-loc-independent-sets-matter-coupling}
A model keeping every Killing-horizon structure and putting local-horizon modes, labelled by null direction, at rapidity period $2\pi c$ satisfies every other present premise, so (R$_{\mathrm{loc}}$) with that labelling is independent of the present premises; Jacobson's argument then gives $G_{\mathrm{cl}} = c/(4\eta) = cG$ at the Killing calibration, so off Killing horizons the clause fixes exactly the equality of the matter coupling with the $G$ of black-hole thermodynamics.
\emph{Source:} Transport Lab claim synthesis.r-loc-independent-sets-matter-coupling, version 3.
\end{fact}

\begin{fact}[The light-ray arrays of the slot-module matter are its Weyl components, sixteen per generation for every CPT halving and every mass term, so a generation registers eight units of ln 2; the particle half is chiral in every halving, so the counted species are not gauge-covariant Dirac species, and with a nonzero Majorana term not Dirac species by mass either]\label{c:synthesis.slot-module-light-ray-arrays-are-weyl-components}
On a horizon's null plane mass terms reach only the constraint, so the arrays of a light ray are the Weyl components; every CPT halving of the slot module has sixteen per generation (eight registered units, for either neutral singlet), and its particle half is chiral in every halving, so the series' matter cannot be grouped into gauge-covariant Dirac species and, with a nonzero Majorana term, is seven Dirac species and two Majorana neutrals per generation.
\emph{Source:} Transport Lab claim synthesis.slot-module-light-ray-arrays-are-weyl-components, version 1.
\end{fact}

\begin{fact}[On a horizon generator the boost is the anticommutator-preserving scaling about the edge, so a stationary gauge-invariant quasi-free state is its own scaling limit there and is the null vacuum of (N$_{+}$) exactly when it is locally the null vacuum at the edge; stationary states thermal on each half at any period, with or without cross-edge correlation, are locally the null vacuum at every other point, so the content of (N$_{+}$) beyond stationarity sits at the bifurcation surface alone]\label{c:synthesis.stationary-horizon-state-is-its-own-edge-scaling-limit}
On every horizon in the scope of (N$_{+}$) the boost is the anticommutator-preserving scaling of a generator about its edge, so a stationary quasi-free state is its own short-distance limit there and satisfies (N$_{+}$) exactly when it is locally the null vacuum at the edge; the stationary states thermal on each half at any period (and at the right period with no cross-edge correlation) are locally the null vacuum at every other point, and zero null energy does not select the vacuum either, so within the stationary class the content of (N$_{+}$) is one temperature-free condition at the bifurcation surface, which on the eternal hole lies off the Painlev\'e-Gullstrand chart.
\emph{Source:} Transport Lab claim synthesis.stationary-horizon-state-is-its-own-edge-scaling-limit, version 2.
\end{fact}

\begin{prediction}[The no-echo prediction, registered]\label{c:scrambling.no-echo-prediction-registered}
Under Postulate (C1), which puts the quantum structure at the stretched singularity, post-merger ringdown is the unmodified Kerr quasinormal spectrum, with no echo train at the delay $8M\ln(M/m_P)$.
\emph{Kill condition:} One statistically confirmed echo train at the Planck-structure delay.
\end{prediction}

\section*{Conclusion}
The thread runs from a dictionary with Einstein's theory, through two routes to his equations, to the quantum state on the horizon that the first route needs.

Each route rests on premises. The horizon route assumes (C$_{\mathrm{loc}}$) and the Clausius relation, and (Loc) is not derived. The action route assumes the field-space premise (FS). Einstein's rule that the whole stress is the source is not given by gauge invariance; it is the trace-free channel, which the horizon route takes through its null fluxes. The temperature comes from (N$_{+}$), whose local form is now a theorem for Hadamard states of the Dirac field. That the series' own state is Hadamard is a further premise, (HD): the state the descent leaves is Hadamard on a Cauchy collar, and propagation carries that to every later field.

Three things stay open. The strength of gravity has no route free of a fitted cut-off. The ground-state clause of (N$_{+}$) is not derived, since local conditions give only the vacuum's family. Nothing in the series yet selects which Hadamard state the world is in.

\end{document}